\documentclass[11pt]{article}

\usepackage[margin=1in]{geometry}
\usepackage{amsmath,amssymb,amsfonts,amsthm}
\usepackage{array}
\usepackage{bm}
\usepackage{multirow}
\usepackage{multicol}
\usepackage{nicefrac}
\usepackage{paralist}
\usepackage{enumitem}
\usepackage{verbatim}
\usepackage{thm-restate}
\usepackage[normalem]{ulem}
\usepackage[compact]{titlesec}

\usepackage{fancybox}
\newenvironment{fminipage}%
  {\begin{Sbox}\begin{minipage}}%
  {\end{minipage}\end{Sbox}\fbox{\TheSbox}}

\newenvironment{algbox}[0]{%
\vskip 0.2in
\noindent
\begin{fminipage}{6.3in}
}{%
\end{fminipage}
\vskip 0.2in
}

\usepackage{xcolor}
\usepackage{cite}
\usepackage{url}
\usepackage{xspace}
\usepackage[mathscr]{euscript}
\usepackage{nameref}
\usepackage[
  linktocpage=true,
  pagebackref=false,
  colorlinks=true,
  linkcolor=blue,
  citecolor=blue,
  urlcolor=blue,
  bookmarks,
  bookmarksopen,
  bookmarksnumbered
]{hyperref}
\usepackage[noabbrev,nameinlink]{cleveref}

% Preserve Computer Modern's existing fallback appearance without font warnings.
\DeclareFontShape{OT1}{cmr}{m}{scit}{<->ssub*cmr/m/sc}{}
\DeclareFontShape{OT1}{cmr}{bx}{sc}{<->ssub*cmr/bx/n}{}

\newtheorem{theorem}{Theorem}[section]
\newtheorem{example}{Example}[section]
\newtheorem{lemma}{Lemma}[section]
\newtheorem{proposition}[lemma]{Proposition}
\newtheorem{corollary}[lemma]{Corollary}
\newtheorem{claim}[lemma]{Claim}
\newtheorem{fact}[lemma]{Fact}
\newtheorem{conjecture}[lemma]{Conjecture}
\newtheorem{definition}[lemma]{Definition}
\newtheorem{problem}{Problem}
\newtheorem{remark}[lemma]{Remark}
\newtheorem{observation}[lemma]{Observation}
\renewcommand*{\theHdefinition}{\theHsection.\arabic{definition}}

\newtheorem*{theorem*}{Theorem}
\newtheorem*{claim*}{Claim}
\newtheorem*{proposition*}{Proposition}
\newtheorem*{lemma*}{Lemma}
\newtheorem*{problem*}{Problem}

\allowdisplaybreaks
\setlength{\parskip}{3pt}

\title{Two-Forest Publisher Maintenance and Transcript-Free Recovery\\
for Incremental Decentralized Data Archival}
\author{}
\date{August 30, 2026}

\begin{document}

\pagenumbering{roman}

\maketitle

\begin{abstract}
We study incremental decentralized data archival for a database of at most
$N$ blocks in $\mathbb F_{2^B}$, initialized with $n_0$ blocks and extended by
one-block Update calls.  For a slack
$\epsilon=\epsilon(\lambda)\in(0,1)$ with polynomially bounded
$1/\epsilon$, we give an explicit Setup--Init--Join--Update--Audit construction
based on tagged dyadic chunks and an additive-evaluation $[Rn,n]$ MDS code,
where $R=\Theta(1/\epsilon)$ and
$R2^{\lceil\log_2N\rceil}\le 2^B$.  Its initialization prefix is frozen in one
forest and the post-initialization suffix is maintained in a second forest.
For every $1\le U\le N-n_0$, node-independent publisher maintenance over the
first $U$ Updates is $O(BU\log N/\epsilon)$ in the simplified block-arithmetic
model, including every public-digest serialization and without a relation
between $n_0$ and $U$; equivalently, the iDDA amortized cost is
$O(B\log N/\epsilon)$ per Update.  Publisher storage is $O(BN/\epsilon)$; the coding
bit-work bound replaces $B$ by $M(B)+B$, where $M(B)$ is the bit cost of field
multiplication.  Setup, Init, and identity-dependent serving are excluded from
this maintenance measure.

For recovery, let $S$ be the advertised node capacity in blocks, let $\mu$ be
the number of pairwise-distinct audited identities, let $\kappa$ be the number
of challenges per sampled chunk, let $\Delta$ bound graph indegree, and let
$\lambda_{\rm vc}$ be the commitment length.  Assume the local
verified-extraction contract for a succinct depth-robust-graph interface,
persistent cache-hidden domain-separated random oracles, quasi-polynomial
vector-commitment binding, polynomial-size public encodings, and polynomially
bounded $\mu,N,\kappa,R$.  Take a public $w=\omega(1)$ satisfying
$w\log_2n_0\ge1$ and
$2(\Delta+1)\kappa w\log_2N\le S$.  Then the same construction has
transcript-free paired recovery provided
$\epsilon\kappa=\omega(\log\lambda)$,
$S=\omega(\kappa\log^2\lambda)$,
$n_0\ge\max\{S,\lambda\}$,
$\lambda_{\rm vc}=\omega(\log\lambda)$, and
$B=\omega(\lambda_{\rm vc}\log N)$.  Its compressor may use the experiment
transcript, whereas its quasi-polynomial extractor receives neither that
transcript nor the original database.  Their only shared randomness is a
finite uniform string, their advice has at most
$B\max\{n-(1-\epsilon)\mu S,0\}$ bits, and
\[
 \Pr[\mathsf{Pass}\wedge(\mathit{DB}_{\rm ext}\ne\mathit{DB}^{*})]
 \le\operatorname{negl}(\lambda).
\]
The theorem concerns recoverability with paired advice and does not assert a
replication bound.
\end{abstract}

\noindent\textbf{Note.} This paper was fully AI generated through a harness created by Richard Peng that discovers open problems and runs them through an automated research pipeline.\par

\clearpage
\tableofcontents
\clearpage

\pagenumbering{arabic}

\section{Introduction}

\label{sec:introduction}

An append-only public database can eventually outgrow the storage available to
an ordinary participant.  Keeping a full replica at every participant then
raises the barrier to entry, whereas merely distributing coded fragments does
not explain which fragments remain available when participants are
untrusted.  Incremental decentralized data archival (iDDA), formulated by
Shi, Silver, and Mu~\cite{ShiSilverMu2026}, brings these issues into one
cryptographic interface.  A publisher initializes a database, storage nodes
join with a fixed space budget, the publisher appends blocks, and periodic
audits test the nodes' evolving shards.  The central recovery requirement is
aggregate: identities that pass their audits should collectively account for
useful database information in proportion to their claimed space, even when
no single identity can store the database.

Several established lines of work capture parts of this problem.  Rabin's
information-dispersal scheme gives threshold reconstruction from coded
fragments~\cite{Rabin1989InformationDispersal}.  Proofs of retrievability
couple successful audits to extraction of an outsourced file, while provable
data possession supplies efficient sampled possession
checks~\cite{JuelsKaliski2007PoR,AtenieseEtAl2007PDP}.  Dynamic proofs of
retrievability support updates and extraction, but provision one server for
the whole file~\cite{CashKupcuWichs2013DynamicPoR,
ShiStefanovPapamanthou2013PracticalDynamicPoR}.  Thus these are nearby
foundations, or full-server restrictions of the archival problem, rather than
results about aggregating many permissionless low-space identities.

Other systems are natural restricted versions of decentralized archival.
Permacoin ties mining rewards to preserving coded portions of a large static
public archive~\cite{MillerJuelsShiParnoKatz2014Permacoin}.  P\'erard et al.,
Kadhe et al., and Raman and Varshney reduce the history stored by an individual
blockchain node and reconstruct from sufficiently many available or honest
participants~\cite{PerardLacanBachyDetchart2018LowStorage,
KadheChungRamchandran2019SeF,RamanVarshney2021Coding}.  Permacoin is static,
while the coded-history guarantees are stated in their respective availability
or corruption models.  None of these works gives recovery proportional to the
space claimed by an arbitrary set of audit-passing pseudonyms.

Our direct same-problem comparison is the iDDA construction of Shi, Silver,
and Mu~\cite{ShiSilverMu2026}.  It provides the five-protocol
Setup--Init--Join--Update--Audit interface and treats recoverability and
replication security as separate objectives.  In its random-oracle
construction, publisher computation is
$B\exp(O(1)/\epsilon)\log N$ per Update, amortized over the prescribed
$N-n_0$ Updates.  This paper addresses the exponential dependence on the
slack in that publisher-maintenance bound.

\paragraph{Our contribution.}
For a database initialized with $n_0$ field elements of $B$ bits and growing
to at most $N$ elements, we construct the five-protocol scheme
$\hyperref[alg:twoforest]{\textsc{TwoForest}}$.  The result permits a varying
slack $\epsilon=\epsilon(\lambda)$ with polynomially bounded reciprocal.  For
every realized prefix of $U$ post-initialization Updates, the scheme maintains
the publisher state and complete public digest using
\[
 O\!\left(\frac{BU\log N}{\epsilon}\right)
\]
node-independent work.  Equivalently, in the amortized terminology of the
direct comparison, publisher maintenance is
$O(B\log N/\epsilon)$ per Update: its dependence on $1/\epsilon$ is linear
rather than exponential.  The prefix guarantee imposes no relation between
$n_0$ and $U$ and is not a worst-case per-call bound.  It excludes one-time
Setup and Init and identity-dependent serving, and it includes every complete
public digest.  Publisher storage is $O(BN/\epsilon)$, and the corresponding
coding bit work over the prefix is
$O((M(B)+B)U\log N/\epsilon)$, where $M(B)$ denotes the bit complexity of
multiplication in $\mathbb F_{2^B}$.

Under the recovery hypotheses stated below, the paired compressor and the
quasi-polynomial extractor
$\hyperref[alg:recovery]{\textsc{TwoForestRecovery}}$ additionally guarantee
that $\mu$ pairwise-distinct identities which all pass their audits can be
combined with at most
\[
 B\max\{n-(1-\epsilon)\mu S,0\}
\]
bits of advice to recover the current $n$-block database, except with
negligible joint error: the probability that all audits pass but the extracted
database is wrong is negligible.  The extractor receives neither the audit
transcript nor the original database.  This is a paired recoverability
statement: it does not establish replication security or a physical-space
lower bound.

We now state the strongest, variable-slack form of the result.  All
asymptotic conditions are for $\lambda\to\infty$, and all logarithms are base
two.

\begin{theorem}[Two-forest publisher and recovery improvement]
\label{thm:main}
Let $\epsilon=\epsilon(\lambda)\in(0,1)$ satisfy
$1/\epsilon\le\lambda^{O(1)}$, let $N>n_0$, and put
\[
 \widehat N=2^{\lceil\log_2N\rceil},\qquad
 k_{\rm rec}=\left\lceil\log_2\frac1\epsilon\right\rceil,
 \qquad \overline\epsilon=2^{-k_{\rm rec}},
\]
\[
 \eta=\frac{\overline\epsilon}{40},\qquad
 \vartheta=1-13\eta,\qquad
 \zeta_0=\frac{\overline\epsilon}{8},\qquad
 R=2^{\lceil\log_2(480\mathrm e/\overline\epsilon)\rceil}.
\]
Assume $R\widehat N\le2^B$, $\lambda_{\rm vc}\le B$, and the
publisher-admissibility conditions of \Cref{def:admissible}.  Then the
five-protocol scheme
$\Pi_{\rm 2F}=\hyperref[alg:twoforest]{\textsc{TwoForest}}$ is honestly correct.
For every $U\in[1,N-n_0]$, the first $U$ post-Init Update calls use
\[
 O\!\left(\frac{BU\log N}{\epsilon}\right)
\]
simplified node-independent publisher work.  This includes each complete
public digest, excludes one-time Setup and Init and identity-dependent serving,
and imposes no horizon condition relating $n_0$ to $U$.  Publisher storage is
$O(BN/\epsilon)$, the coding bit work is
$O((M(B)+B)U\log N/\epsilon)$, and a full-codeword Merkle tree has depth
$O(\log N+\log(1/\epsilon))$.  Thus, in the iDDA amortized accounting, the
simplified cost is $O(B\log N/\epsilon)$ per Update; the claim is not a
worst-case per-call bound.

Suppose in addition that the family is recovery-admissible in the sense of
\Cref{def:admissible}, that the local verified-extraction contract of
\Cref{def:local-extraction} holds uniformly at every legal current length, and
that
\[
 \epsilon\kappa=\omega(\log\lambda),\quad
 1\le\mu\le\lambda^{O(1)},\quad
 S=\omega(\kappa\log^2\lambda),\quad
 n_0\ge\max\{S,\lambda\},
\]
\[
 w=\omega(1),\qquad
 \lambda_{\rm vc}=\omega(\log\lambda),\qquad
 B=\omega(\lambda_{\rm vc}\log N),
\]
and that the $\mu$ challenge identities are pairwise distinct.  Then the two
entry points of
\hyperref[alg:recovery]{\textsc{TwoForestRecovery}} are total and have the
signatures
\[
 \mathcal C^{\mathcal A}(1^\lambda,S,\mu,\mathit{tr};\rho)
      \longrightarrow\mathit{DB}_{\rm short},
\]
\[
 \mathcal E^{\mathcal A_2(\mathit{st}_{\mathcal A})}
 (1^\lambda,n,S,\mu,\varphi_n,\mathit{DB}_{\rm short};\rho)
      \longrightarrow\mathit{DB}_{\rm ext}\ \text{or }\bot.
\]
Here $\rho$ is finite uniform shared randomness independent of the receiver
experiment.  The extractor receives neither $\mathit{tr}$ nor
$\mathit{DB}^{*}$, may face arbitrary adaptive dependence of later audits on
earlier audits, and runs in time quasi-polynomial in
$(\lambda,t_{\mathcal A})$.  For advice emitted by the paired compressor in the
same public, continuation, persistent-oracle, and shared-coin context,
\[
 |\mathit{DB}_{\rm short}|
 \le B\max\{n-(1-\epsilon)\mu S,0\},
\]
and
\[
\Pr[\mathsf{Pass}\wedge
       (\mathit{DB}_{\rm ext}\ne\mathit{DB}^{*})]
 \le\operatorname{negl}(\lambda).
\]
At these parameters an honest node state has $O(BS)$ bits.  If an Update
refreshes $d$ honest participants, its separately charged identity-dependent
output is $O(dBS)$ bits and its complete explicit output, including the one
public digest, is $O((d+1)BS)$ bits.  Here the publisher hierarchy is updated
in place and is not reserialized as output.
No correctness assertion is made for unpaired advice, and no replication
probability conclusion is part of the theorem.
\end{theorem}

\paragraph{Proof ideas and scope.}
The publisher construction freezes the initialized prefix in one tagged
dyadic forest and maintains only the appended suffix in a second binary-counter
forest.  An additive-evaluation MDS code with redundancy
$R=\Theta(1/\epsilon)$ makes equal-size encoded chunks recursively mergeable;
the binary-counter accounting makes each appended block participate in only
$O(\log N)$ merge levels, yielding $O(BR\log N)$ average work.  For recovery,
accepted audits yield candidate codeword symbols, an accepting-spine argument
handles adaptive audit dependence without giving the transcript to the
extractor, and collision expansion followed by error-and-erasure decoding
completes each chunk.  These ingredients prove the stated paired-advice bound.
The construction of Shi, Silver, and Mu also proves a replication-security
result~\cite{ShiSilverMu2026}, whereas
\Cref{thm:main} assumes distinct challenged identities and proves
recoverability only; consequently, the publisher improvement is not an
overall strengthening of their theorem.

\paragraph{Organization.}
\Cref{sec:preliminaries} fixes the computational model, the recovery
experiment, and the accounting conventions.  \Cref{sec:overview} explains the
publisher and recovery arguments.  \Cref{sec:main-proof} constructs the
five-protocol scheme and proves \Cref{thm:main}.

\section{Preliminaries}
\label{sec:preliminaries}

\paragraph{Notation.}
All logarithms are base two.  For integers $a\le b$, write
$[a,b]=\{a,a+1,\ldots,b\}$ and
$[a,b)=\{a,a+1,\ldots,b-1\}$; in particular, $[m]=[0,m)$.
The symbol $\Vert$ denotes concatenation, and $\langle x\rangle$ denotes the
canonical prefix-free binary serialization of $x$; its length is
$|\langle x\rangle|$.  If $A$ is an array on a finite ordered set $\Omega$
and $Q=(q_1,\ldots,q_r)$ is an index tuple, then
$A[Q]=(A[q_1],\ldots,A[q_r])$, retaining order and multiplicity.  If
$Q\subseteq\Omega$ is a set, then $A[Q]=A|_Q$ uses the order induced by
$\Omega$.

\paragraph{Databases and computation.}
A database of length $n$ is
$\mathit{DB}_n=(b_0,\ldots,b_{n-1})\in\mathbb F^n$, where
$\mathbb F=\mathbb F_{2^B}$ is represented in a fixed polynomial basis.
Thus field addition is bitwise XOR, denoted $\oplus$.  Initialization installs
$\mathit{DB}_{n_0}$, and after $u$ one-block Update calls the current length is
$n=n_0+u\le N$.  The quantity $M(B)$ is the bit complexity of one
multiplication in $\mathbb F$.  The commitment and Merkle-root length is
$\lambda_{\rm vc}$ bits.  A function is negligible if it is
$\lambda^{-\omega(1)}$.  All asymptotic conditions are for
$\lambda\to\infty$, and ``quasi-polynomial'' means
$2^{(\log\lambda)^{O(1)}}$ times a polynomial in the other displayed input
sizes.

\paragraph{Partial words and MDS decoding.}
For a finite ordered set $\Omega$, a partial word is a map
$Y:\Omega\to\mathbb F\cup\{\bot\}$ with
\[
 \operatorname{supp}(Y)=\{q\in\Omega:Y[q]\ne\bot\}.
\]
Relative to a word $C\in\mathbb F^\Omega$, its error count is
\[
 e_C(Y)=|\{q\in\operatorname{supp}(Y):Y[q]\ne C[q]\}|.
\]
An $[M,k]$ MDS code has the property that every $k$ coordinates determine its
message.  We use the standard error-and-erasure criterion: if $Y$ is a partial
received word relative to a transmitted codeword $C$ and
$P=\operatorname{supp}(Y)$, then $C$ is the unique codeword consistent with at
most $e_C(Y)$ errors whenever
\[
 |P|-2e_C(Y)\ge k.
\tag{EE}\label{eq:prelim-error-erasure}
\]
We also use that a nonzero polynomial of degree at most $d$ over a field has
at most $d$ roots.

\paragraph{Typed objects and persistent random functions.}
A committed array is identified by its type $\gamma$, length $m$, and ordered
contents $A$.  The type includes the object's role, full interval tag, and all
applicable representation, identity, and context fields.  With public
reference string $\mathit{crs}$, the commitment digest algorithm produces a
$\lambda_{\rm vc}$-bit root and opening state; an opening names an aligned
coordinate set $Q$, the values $A[Q]$, and a proof.  We assume correctness and
binding against quasi-polynomial algorithms for this complete typed instance.
In particular, except with the binding advantage, two accepted openings of
one typed root cannot assign different values to the same coordinate.

Random oracles are persistent fixed random functions.  Equivalently, in a
lazy implementation each fresh typed input receives an independent uniform
answer, and prefix-free domain tags name mutually independent families.
Cache membership and query order are not visible to a machine.  For a fixed
typed audit context $\Xi$, let $\Gamma_{\mathcal B}$ be the space of complete
configurations of a continuation $\mathcal B$, including its random tape,
work tapes, and communication state, but not the ambient oracle tables.
Restoring $\sigma\in\Gamma_{\mathcal B}$ restores only the machine: it neither
erases nor reprograms any oracle.  Once $\Xi$ and the persistent functions are
fixed, the transition maps of the continuation are deterministic.

\begin{definition}[Distinct-identity receiver experiment and paired recovery]
\label{def:receiver-experiment}
Fix a deterministic non-uniform polynomial-time adversary
$\mathcal A=(\mathcal A_1,\mathcal A_2)$, and let
$t_{\mathcal A}(\lambda)$ be a polynomial upper bound on either phase and on
each resumed continuation segment.  In the first phase of the iDDA receiver
experiment, the challenger and $\mathcal A_1$ execute the public setup,
initialization, and adaptive Join and Update interactions of the scheme.  At
an adversarially chosen stopping length $n$, this phase fixes the original
database $\mathit{DB}^{*}$, its public digest $\varphi_n$, the continuation
state $\mathit{st}_{\mathcal A}$, and a transcript $\mathit{tr}$ containing
the challenger coins and messages needed to replay the phase.  The second
phase runs
$\mathcal B=\mathcal A_2(\mathit{st}_{\mathcal A})$ through $\mu$ sequential
Audit interactions at this same public state.  We write
\[
 (b,\varphi_n,\mathit{DB}^{*},\mathit{st}_{\mathcal A},\mathit{tr})
 \leftarrow
 \mathsf{RecvExpt}^{\mathcal A}(1^\lambda,S,\mu)
\]
for the resulting tuple.

In each Audit, the continuation fixes its identity and offline message before
the online challenge.  Conditional on those values, every randomly challenged
active object receives an independent exact-uniform $\kappa$-tuple in its
declared finite range; a full-opening object has a deterministic challenge.
Later identities, messages, and checkpoints may depend on all earlier Audit
interactions.  A repeated identity makes the experiment reject, and $b=1$
exactly on the event $\mathsf{Pass}$ that all $\mu$ distinct-identity Audits
accept.

The recovery probability space first samples this experiment, including its
setup coins, persistent random functions, and Audit challenges, then samples
the finite uniform string $\rho$ independently and runs the paired compressor
and extractor displayed in \Cref{thm:main}.  Thus every probability below is
over these choices.  The compressor may use $\mathit{tr}$ to replay
$\mathcal A_1$.  The extractor receives neither $\mathit{tr}$ nor
$\mathit{DB}^{*}$, but only $\mathcal B$ and its displayed public inputs and
advice.  Both algorithms use the experiment's ambient public parameters and
persistent functions.  Inputs are \emph{paired} when the advice was produced
by the compressor in this same public, continuation, oracle, and shared-coin
context; transcript replay must reconstruct the corresponding digest,
continuation state, and prior persistent-oracle context.
\end{definition}

\paragraph{Cost conventions.}
Publisher maintenance comprises construction and maintenance of encoded
hierarchy objects, their Merkle trees, and every serialized public digest.  In the
simplified block-arithmetic model, one field operation and one hash on
$O(B+\lambda_{\rm vc})$ input bits each cost $O(B+\lambda_{\rm vc})$ work;
under $\lambda_{\rm vc}\le B$, displayed coding bounds use a factor $B$.
Serialization is charged by its actual bit length.  In the genuine coding
bit model, the field-operation factor is $M(B)+B$.  Post-Init publisher
maintenance excludes one-time Setup and Init and accounts for
identity-dependent Join serving separately.

\section{Overview}
\label{sec:overview}

The proof has a publisher-maintenance part and a paired-recovery part.  The
first combines a prefix-separated binary-counter representation with the
recursively mergeable MDS code
\hyperref[alg:additive]{\textsc{AdditiveEC}}.  The second uses the exact finite
coin sampler \hyperref[alg:challenge-sample]{\textsc{TaggedChallengeSample}},
the accepting-spine rewind procedure
\hyperref[alg:adaptive-raw]{\textsc{AdaptiveRawExtract}}, and MDS
error-and-erasure completion in
\hyperref[alg:recovery]{\textsc{TwoForestRecovery}}.  The explicit local
verified-extraction contract in \Cref{def:local-extraction} supplies many
usable occurrences; unconditional collision expansion converts them into
enough distinct coordinates despite adaptive spine selection; and no-index
MDS completion recovers each chunk.  The local contract is a hypothesis,
whereas this coupling, expansion, and completion chain is established in the
proof.  To obtain the strongest, variable-slack statement, we round the target
slack down to a public dyadic value and set
\[
 \overline\epsilon=2^{-\lceil\log_2(1/\epsilon)\rceil},\qquad
 \eta=\frac{\overline\epsilon}{40},\qquad
 \vartheta=1-13\eta,\qquad
 R=2^{\lceil\log_2(480\mathrm e/\overline\epsilon)\rceil}.
\]
Thus $\epsilon/2<\overline\epsilon\le\epsilon$ and
$R=\Theta(1/\epsilon)$.  The same rounding makes
$\zeta_0=\overline\epsilon/8$ use only
$O(1+\log(1/\epsilon))$ header bits, which the field-size condition lets the
publisher bound absorb.

\paragraph{Separating the initialized prefix from later carries.}
At length $n=n_0+u$, \textsc{TwoForest} represents the nonzero binary digits
of $n_0$ and $u$ as two independently tagged dyadic forests.  The first forest
partitions the initialized interval $[0,n_0)$ and is frozen after Init; the
second starts empty and partitions the appended interval $[n_0,n_0+u)$.
Their $O(\log N)$ chunks, tagged by the forest, global starting point, and
level, partition the database and are joined in chronological order by
$\mathsf{TwoForestJoin}$ (\Cref{def:chunks,lem:partition}).  This separation
addresses the publisher-side obstacle: in a single initialized binary counter,
a post-Init carry can propagate through a large portion of the old prefix, so
its cost cannot be bounded in terms of the number $U$ of subsequent Updates
alone.  Here every such carry is confined to the initially empty suffix
forest, with no condition relating $n_0$ and $U$.

The code makes each suffix carry cheap.  Its recursive additive-polynomial
construction, based on $\psi_t(X)=X^2+c_tX$, evaluates on nested spaces $V_t$
to give an $[R2^t,2^t]$ MDS word, whose first $2^{t+1}$ coordinates also form a
redundancy-two MDS word.  Its essential identity combines two equal-size child
words coordinatewise as
\[
 \mathsf{Merge}_t(C^{(0)},C^{(1)})[z]
 =C^{(0)}[\psi_t(z)]+zC^{(1)}[\psi_t(z)],
\]
so a size-$2^t$ parent costs only $O(R2^t)$ field operations
(\Cref{lem:additive-code}).  Every chunk caches its raw message, the
redundancy-two restriction, and the full redundancy-$R$ word, together with
their typed trees.  This is needed because the audit class of a persistent
chunk can change as $n$ changes; Update can then select a cached root instead
of re-encoding that chunk.

If level $\ell$ is created during a suffix carry, it costs
$O(BR2^\ell)$ in the simplified model and is created at most
$\lfloor U/2^\ell\rfloor$ times in the first $U$ Updates.  Hence
\[
\sum_{\ell=0}^{\lfloor\log_2U\rfloor}O(BR2^\ell)
       \left\lfloor\frac{U}{2^\ell}\right\rfloor
 =O(BRU\log N)=O\!\left(\frac{BU\log N}{\epsilon}\right).
\]
Serializing all $U$ public digests adds
$O(BU\log N+U|\langle\zeta_0\rangle|)$ work; the dyadic encoding and field-size
condition absorb this term into the displayed bound.
The implementation in \hyperref[alg:twoforest]{\textsc{TwoForest}} runs
Setup, Init, Join, Update, and Audit on these objects: Setup prepares the code
and typed namespaces, Init builds and freezes the prefix forest, Update changes
only the suffix forest, Join materializes an identity's authenticated samples
and labels, and Audit checks the typed offline and online openings.  The same
simplified node-independent post-Init metric treats one-time Setup and Init and
identity-dependent Join serving separately.  It gives $O(BN/\epsilon)$
publisher storage and replaces $B$ by $M(B)+B$ for coding bit work; a full
codeword tree has depth $O(\log N+\log(1/\epsilon))$.  It is an aggregate,
hence amortized, Update bound rather than a worst-case bound for an individual
carry.

\paragraph{From accepted audits to candidate codeword symbols.}
For a sampled chunk $c$ of length $n_c$, an identity stores
\[
 s_c=\left\lceil\max\left\{\frac Sn n_c,s_{\rm lb}(n)\right\}\right\rceil
\]
positions sampled independently with replacement from its length-$Rn_c$ word,
where
$s_{\rm lb}(n)$ is the public lower cutoff in
\Cref{eq:lower-sampling}.  The cutoff ensures that a sampled chunk has at
least $\kappa$ occurrences for the extraction estimates.  Together, the
proportional term, this cutoff, the recovery-validity bounds, and the
$O(\log N)$-chunk decomposition keep total honest-node data at $O(S)$.
Smaller mini chunks use their whole redundancy-two word, and tiny chunks use
their raw word, avoiding a sampling argument in ranges too small to support
it.  For each identity, the sampled data are masked into the final quarter of
a depth-robust-graph label array.  Writing $\widetilde h_{c,i}$ for the final
verified label array accumulated by the rewind procedure below, an accepted
online opening of a terminal label and all its parents yields the candidate
pair
\[
 \left(g_{c,\mathit{id}_i}(j),\;
 \widetilde h_{c,i}[q]\mathbin\oplus
 H_c((\varphi_n,\mathit{id}_i),q,
       \widetilde h_{c,i}[\operatorname{Par}_c(q)])\right),
 \qquad q=3s_c+j.
\]
Typed commitment binding makes accepted openings consistent at an opened
coordinate, but the candidate need not be correct: usability is an analysis
notion that the algorithm neither knows nor tests.  The local extraction
contract asserts that, on its good event, at least $(1-4\eta)s_c$ occurrences
per selected sampled identity are usable; it gives the analogous
redundancy-two and full-raw conclusions for mini and tiny chunks.

\paragraph{A transcript-free accepting spine.}
The recovery algorithms must agree on those selected identities and rewind
points even though later audits may depend on earlier interactions and the
extractor receives neither the receiver transcript nor the original database.
\textsc{TaggedChallengeSample} first parses the finite shared uniform string
$\rho$ by bounded rejection.  This is necessary because the challenge ranges
need not have power-of-two size; conditional on the negligible parser-failure
event not occurring, all named challenges are exactly uniform, mutually
independent, and independent of the receiver experiment.

Given the public audit context, the continuation, the public allocation, and
$\rho$, \textsc{AdaptiveRawExtract} tries $K_{\rm sp}$ complete audit paths
driven by independent challenge tables, conditional on the fixed public
context and persistent random functions.  It selects the least path on which
all $\mu$ distinct-identity audits accept and saves the machine configuration
immediately after each identity submission and before its offline message.
If a fresh path accepts with conditional probability $p$, its coupling to the
received accepting execution is controlled by
\[
 p(1-p)^{K_{\rm sp}}\le\frac1{K_{\rm sp}+1}.
\]
The choice of $K_{\rm sp}$ in \Cref{eq:spine-count} makes this loss negligible
while remaining quasi-polynomial.  A checkpoint wrapper then repeatedly
restores each saved machine configuration for the $T_{\rm ext}$ raw trials,
without resetting the persistent random functions, and aggregates only values
from trials whose offline and online predicates both accept.  Its output is a
partial candidate array for every sampled or mini chunk and a partial raw array
for every tiny chunk.  Once the supplied continuation state and public context
are fixed, this computation uses no transcript.  Its deterministic least-spine
rule and last-write recurrence therefore give paired callers with the same
persistent functions and $\rho$ exactly the same arrays in quasi-polynomial
time (\Cref{lem:spine}).

\paragraph{Removing postselection bias.}
The central recovery difficulty remains that the identities on the accepting
spine can themselves depend on answers of the shard random function $G_c$.
One therefore cannot condition on spine success and then apply an ordinary
independent occupancy bound to their sampled positions.  The proof instead
exposes the adaptive set of first arguments reached by the experiment, replay,
and spine search, pads it to a deterministic size, and defers the local rewind
queries until after this exposure.  It can then apply the deletion-resilient
collision lemma unconditionally.  Since $R\ge12\mathrm e/\eta$, retaining any
$(1-4\eta)s_c$ usable occurrences from each of the selected $r_c$ identities
still leaves at least $(1-5\eta)r_cs_c$ distinct coordinates
(\Cref{lem:collision,lem:deferred}).

The public allocation is chosen to turn that expansion into precisely the
decoding margin that is needed:
\[
 r_c=\min\left\{\mu,
       \left\lceil\frac{n_c}{\vartheta s_c}\right\rceil\right\},\qquad
 c_c=\min\{n_c,\lceil\vartheta\mu s_c\rceil\},\qquad
 d_c=n_c-c_c.
\]
Let $Y_c$ keep the first populated candidate at each coordinate, let $P_c$ be
its support, and let $e_c$ count its unknown errors relative to the authentic
codeword.  An unusable occurrence can cause at most one bad first write.
Combining the distinct-coordinate bound with the at most $4\eta r_cs_c$
unusable occurrences gives the pivotal estimate
\[
 |P_c|-2e_c\ge(1-13\eta)r_cs_c
 =\vartheta r_cs_c,
\]
and integrality together with the definition of $r_c$ upgrades this to
$|P_c|-2e_c\ge c_c$.
This is why collision expansion must survive both adaptive selection and
deletion: support alone is insufficient when the extracted partial word can
contain errors that the extractor cannot recognize.

\paragraph{No-index completion and final assembly.}
For every sampled chunk, both paired callers choose $I_c$ to be the first
$d_c$ public evaluation points outside $P_c$.  The compressor, which may use
the transcript to recover the original database, sends only the authentic
codeword values at $I_c$ in canonical order.  It sends neither $I_c$ nor any
support metadata.  The transcript-free extractor recomputes the same set from
the shared raw record, fills those values into $Y_c$, and obtains
\[
 (|P_c|+d_c)-2e_c\ge c_c+(n_c-c_c)=n_c,
\]
which is exactly the MDS error-and-erasure criterion.  Mini chunks meet the
same criterion in their $[2n_c,n_c]$ restriction, tiny chunks are recovered by
their full raw openings, and $\mathsf{TwoForestJoin}$ reassembles the decoded
messages in database order.

Finally, $s_c\ge(S/n)n_c$ and the definition of $c_c$ imply
\[
 B\sum_{c\ {\rm sampled}}d_c
 \le B\max\{n-(1-\epsilon)\mu S,0\}.
\]
The sampler, spine, local-extraction, binding, and unconditional-expansion
losses are jointly negligible under the hypotheses of \Cref{thm:main}, which
gives
$\Pr[\mathsf{Pass}\wedge(\mathit{DB}_{\rm ext}\ne\mathit{DB}^{*})]
\le\operatorname{negl}(\lambda)$ for the quasi-polynomial extractor.  This
conclusion requires the stated pairwise-distinct challenge identities and
paired advice in the same public, continuation, persistent-oracle, and
shared-coin context.  It is a recoverability theorem; no replication
probability conclusion is used or proved.

\section{Proof of the Main Theorem}
\label{sec:main-proof}

\subsection{Parameters and the two-forest decomposition}

We first isolate the public hypotheses.  The local graph-extraction statement
used later is recorded separately in \Cref{def:local-extraction}; it is a
hypothesis of the recovery clause, not of publisher accounting.

\begin{definition}[Admissible public parameters and recovery interface]
\label{def:admissible}
Fix an absolute constant $c_\Delta>0$.
A publisher-admissible family consists, at each security parameter $\lambda$,
of integers $B,\allowbreak S,\allowbreak n_0,\allowbreak N,\allowbreak
\kappa,\allowbreak\Delta$, a public integer $w=w(\lambda)$, a
$\lambda_{\rm vc}$-bit typed Merkle vector commitment, and the displayed
values $R,\zeta_0$ from \Cref{thm:main}.  It satisfies
\[
 2\le\lambda\le n_0<N\le\lambda^{O(1)},\qquad
 1\le S\le n_0,\qquad 1\le\kappa,w\le\lambda^{O(1)},
\]
\[
 1\le\Delta\le c_\Delta\log_2(2N),\qquad
 w\log_2n_0\ge1,\qquad
 2(\Delta+1)\kappa w\log_2N\le S,
\tag{1}\label{eq:recovery-validity}
\]
as well as $R\widehat N\le2^B$, $1\le\lambda_{\rm vc}\le B$, and
$1\le|\langle\zeta_0\rangle|\le BS$.  The public
directed-graph interface, on a length $m\le8N$, a domain tag, and a vertex
$q\in[m]$, returns in polynomial time an increasing list of at most $\Delta$
parents, all smaller than $q$.  For each tag these lists define the promised
depth-robust graph.  The graph description $\mathit{pp}_{\rm DRG}$ has
$O(\lambda+\log N)$ bits.
For purposes of the recovery theorem, the semantic content of this graph
promise is exactly the uniform verified-extraction statement in
\Cref{def:local-extraction}; publisher correctness uses only acyclicity and the
displayed indegree and running-time bounds.
The Merkle leaf hash binds a prefix-free encoding of the object type, full
interval tag, array length, coordinate, and value.  Public setup creates the
field tables and independently domain-separated graph and random-oracle
namespaces.

Call a publisher-admissible family \emph{recovery-admissible} if, in addition,
$N,R,\kappa,\mu\le\lambda^{O(1)}$, every serialized digest has polynomial
length, and the complete public recovery input has length
$\operatorname{poly}(\lambda,t_{\mathcal A})$.  Its public Audit interaction
is represented by total deterministic transition maps
$\mathsf{NextId}$, $\mathsf{Offline}$, and $\mathsf{Online}$ once the persistent
random functions are fixed.  A call, including parsing, communication, graph
operations, oracle calls, and verification over all active chunks, has bit
cost at most $C_{\rm tr}$, where
$L_{\varphi_n}:=|\varphi_n|$ is the actual serialized current-digest length and
\[
 C_{\rm tr}=t_{\mathcal A}(\lambda)
   +\operatorname{poly}(\lambda,N,R,B,\kappa,L_{\varphi_n})
\tag{2}\label{eq:transition-cost}
\]
Malformed or overlong records have total rejecting branches.
The typed vector commitment is correct and binding against quasi-polynomial
algorithms, and all random-oracle families are persistent, cache hidden, and
mutually domain separated.  These input-size conditions ensure efficiency;
they do not give either recovery algorithm additional information.
\end{definition}

At length $n=n_0+u$, the prefix and suffix binary expansions give the following
tagged chunks.  This dyadic organization is the binary-counter pattern behind
static-to-dynamic transformations~\cite{BentleySaxe1980StaticDynamic}, applied
separately on the two sides of the initialization boundary.

\begin{definition}[Prefix-separated tagged chunks]
\label{def:chunks}
For $0\le u\le N-n_0$, let
\[
 H(v)=\{\ell:\text{bit }\ell\text{ of }v\text{ is }1\},\qquad
 \mathcal K(u)=\{(0,\ell):\ell\in H(n_0)\}
   \cup\{(1,\ell):\ell\in H(u)\}.
\]
For $c=(f,\ell)\in\mathcal K(u)$ put
$f_c=f$, $\ell_c=\ell$, $n_c=2^{\ell_c}$, and
\[
 a_{(0,\ell)}=\sum_{\substack{j\in H(n_0)\\j>\ell}}2^j,
 \qquad
 a_{(1,\ell)}=n_0+
   \sum_{\substack{j\in H(u)\\j>\ell}}2^j.
\]
For $\mathit{DB}=(b_0,\ldots,b_{n-1})$, define
\[
 x_{n,c}=(b_{a_c},\ldots,b_{a_c+n_c-1}),
 \qquad \iota_{n,c}=(f,a_c,\ell,n_c).
\tag{3}\label{eq:chunk-data}
\]
Chunks are ordered first by $f$ and then by decreasing $\ell$.  Denote the
resulting split and concatenation maps by $\mathsf{TwoForestSplit}_{n_0,u}$ and
$\mathsf{TwoForestJoin}_{n_0,u}$.
\end{definition}

The binary intervals in this definition have the following elementary but
essential reconstruction property.

\begin{lemma}[Tagged chunks partition the database]
\label{lem:partition}
For every $u\in[0,N-n_0]$,
\[
 |\mathcal K(u)|\le2(1+\lfloor\log_2N\rfloor),\qquad
 \sum_{c\in\mathcal K(u)}n_c=n_0+u,
\]
and
\[
 \mathsf{TwoForestJoin}_{n_0,u}
 \bigl(\mathsf{TwoForestSplit}_{n_0,u}(\mathit{DB})\bigr)=\mathit{DB}.
\]
\end{lemma}

\begin{proof}
Within either forest, the decreasing nonzero binary digits describe disjoint
consecutive intervals whose lengths sum to the represented integer.  The first
forest partitions $[0,n_0)$ and the second partitions $[n_0,n_0+u)$.  Each
binary expansion has at most $1+\lfloor\log_2N\rfloor$ nonzero digits.  The
three assertions follow.
\end{proof}

The next ingredient lets the binary carries update a codeword without
re-encoding all of its descendants.  It is a polynomial-evaluation MDS code in
the Reed--Solomon tradition~\cite{ReedSolomon1960PolynomialCodes}, arranged so
that equal-size words admit a linear-work merge.

\subsection{A recursively mergeable MDS code}

Let $L_{\max}=\lceil\log_2N\rceil$.  Fix an ordered
$\mathbb F_2$-basis of $\mathbb F$.
Since $R\widehat N\le2^B$ and $R$ is a power of two, let
$V_{L_{\max}}\subseteq\mathbb F$ be the span of the first
$\log_2(R\widehat N)$ basis vectors.  For $t=L_{\max},\ldots,1$, take $c_t$
to be the first nonzero vector in the current row-reduced basis of $V_t$ and
set
\[
 \psi_t(X)=X^2+c_tX,\qquad V_{t-1}=\psi_t(V_t).
\]
The restriction of $\psi_t$ to $V_t$ is linear with kernel $\{0,c_t\}$, so
$|V_t|=R2^t$.  Compute the canonical row-reduced basis of $V_{t-1}$ and the
canonical row-reduced linear section $\sigma_t:V_{t-1}\to V_t$.  Order $V_t$
by placing $\sigma_t(y),\sigma_t(y)+c_t$ consecutively for each point $y$ in
the fixed order of $V_{t-1}$, starting with lexicographic field-coordinate
order on $V_0$.  Write $z_{t,q}$ for the $q$th point.

For $x\in\mathbb F^{2^t}$ define a polynomial recursively by
\[
 P_{(x_0)}^{(0)}(X)=x_0,\qquad
 P_x^{(t)}(X)=P_{x^{(0)}}^{(t-1)}(\psi_t(X))
       +X P_{x^{(1)}}^{(t-1)}(\psi_t(X)),
\tag{4}\label{eq:additive-polynomial}
\]
where $x=x^{(0)}\Vert x^{(1)}$.  Define
\[
 \mathsf{Enc}_{2^t}(x)[z]=P_x^{(t)}(z)\quad(z\in V_t),
\]
and let $W_t$ be the first $2^{t+1}$ points of $V_t$.  For arrays
$C^{(0)},C^{(1)}\in\mathbb F^{V_{t-1}}$, define
\[
 \mathsf{Merge}_t(C^{(0)},C^{(1)})[z]
 =C^{(0)}[\psi_t(z)]+zC^{(1)}[\psi_t(z)].
\tag{5}\label{eq:merge}
\]
Public setup stores each $z$ together with the position of $\psi_t(z)$ in
$V_{t-1}$.  These tables contain $O(R\widehat N)$ field elements and indices.

\begin{algbox}
\textbf{\uline{\textsc{AdditiveEC}}}: Set up, encode, merge, or decode an
additive-evaluation word.
\phantomsection\label{alg:additive}

\uline{Input}: A mode
$m\in\{\mathsf{setup},\mathsf{encode},\mathsf{merge},\mathsf{decode}\}$ and a
mode-specific payload: $(B,N,R)$; $(t,x)$; $(t,C^{(0)},C^{(1)})$; or a partial
word $(t,Y)$, respectively.  The public field representation and, outside
setup mode, the tables above are ambient.

\uline{Output}: Public encoding tables; $\mathsf{Enc}_{2^t}(x)$;
$\mathsf{Merge}_t(C^{(0)},C^{(1)})$; or the unique message within the stated
error-and-erasure radius or $\bot$, respectively.

\uline{Guarantee}: In decode mode, if $P=\{z:Y[z]\ne\bot\}$ and $Y|_P$ has
$e$ errors relative to a transmitted word, the transmitted message is returned
whenever the criterion in \Cref{eq:prelim-error-erasure} holds, namely
$|P|-2e\ge2^t$.

\begin{enumerate}[leftmargin=*,label=\arabic*.]
 \item In setup mode, construct the spaces, sections, orders, and lookup tables
 described above and output them.
 \item In encode mode, recursively encode the two halves of $x$ and combine
 them coordinatewise by \Cref{eq:merge}; at $t=0$, evaluate the constant
 polynomial on $V_0$.
 \item In merge mode, scan the level-$t$ table and evaluate \Cref{eq:merge}
 once at each output coordinate.
 \item In decode mode, output $\bot$ if $|P|<2^t$.  Otherwise set
 $H=\lfloor(|P|-2^t)/2\rfloor$.  For $h=0,\ldots,H$, solve the
 Berlekamp--Welch equations
 $Q(z)=Y[z]E(z)$ on $P$, where $E$ is monic of degree $h$ and
 $\deg Q<2^t+h$.  Test every quotient $Q/E$ of degree less than $2^t$ and
 apply the inverse of the fixed coefficient map $x\mapsto P_x^{(t)}$ to the
 unique candidate with at most $H$ disagreements and output the resulting
 message; output $\bot$ if no such candidate exists.
\end{enumerate}
\end{algbox}

The next lemma proves both the algebraic guarantees and the resource bounds of
the displayed implementation.

\begin{lemma}[Linear-work additive MDS coding]
\label{lem:additive-code}
The map $\mathsf{Enc}_n$, for $n=2^t$, is an $[Rn,n]$ MDS code, and its
restriction to $W_t$ is an $[2n,n]$ MDS code.  The algorithm
\hyperref[alg:additive]{\textsc{AdditiveEC}} is correct.  Setup uses
$O(L_{\max}B^3)$ binary-linear-algebra work plus $O(R\widehat N)$ field
operations and table writes; Merge uses $O(Rn)$ field operations; Encode uses
$O(Rn(1+\log n))$ field operations; and the displayed direct decoder uses at
most $O((Rn)^4+n^3)$ field operations.
\end{lemma}

\begin{proof}
Induction in \Cref{eq:additive-polynomial} gives $\deg P_x^{(t)}<2^t$.  To
prove injectivity, suppose the polynomial is zero and abbreviate its two child
polynomials by $A$ and $D$.  The two preimages of $y\in V_{t-1}$ are $z$ and
$z+c_t$, and characteristic two gives
\[
 0=P_x^{(t)}(z)+P_x^{(t)}(z+c_t)=c_tD(y).
\]
Thus $D$ vanishes at $|V_{t-1}|=R2^{t-1}\ge2^{t-1}$ points while
$\deg D<2^{t-1}$, so $D=0$; then $A=0$.  Induction recovers $x=0$.
Consequently the messages parameterize all polynomials of degree less than
$n$, and any $n$ distinct evaluations determine the message.  This proves both
MDS assertions.

\Cref{eq:additive-polynomial} is exactly the merge rule in
\Cref{eq:merge}.  One
multiplication and one addition per coordinate give the merge bound, and
$T(n)=2T(n/2)+O(Rn)$ gives the encoding bound.  If two degree-less-than-$n$
polynomials each disagree with $Y$ at at most
$H=\lfloor(|P|-n)/2\rfloor$ positions, they agree at at least $n$ points and
hence are equal.  If the actual number of errors is $e$ and $|P|-2e\ge n$, the
error locator of degree $e$ appears in the search; the usual root count forces
$Q=EP_x^{(t)}$, so division returns the transmitted polynomial.  Direct
Gaussian elimination gives the conservative decoding bound.  Binary row
reduction constructs every image and section, and the geometric sum of table
sizes is $O(R\widehat N)$, proving the setup bounds.
\end{proof}

We now use this code inside the exact tagged state maintained by the
publisher.

\subsection{The five-protocol two-forest construction}

At a current length $n$, put $L(n)=\lfloor\log_2n\rfloor$ and
\[
 s_{\rm lb}(n)=\left\lceil
  \frac{S2^{L(n)}}{n w\log_2n}\right\rceil.
\tag{6}\label{eq:lower-sampling}
\]
A chunk is \emph{tiny} if $n_c\le\kappa$, \emph{mini} if
$\kappa<n_c\le s_{\rm lb}(n)$, \emph{small} if
$s_{\rm lb}(n)<n_c\le s_{\rm lb}(n)N/S$, and \emph{large} otherwise.
Small and large chunks are called \emph{sampled}.  For a sampled chunk set
\[
 s_c=\left\lceil\max\left\{\frac Sn n_c,s_{\rm lb}(n)\right\}\right\rceil,
 \qquad \tau_c=s_c,\qquad m_c=Rn_c.
\tag{7}\label{eq:sample-size}
\]
For a mini chunk put $(\tau_c,m_c)=(2n_c,2n_c)$, and for a tiny chunk put
$m_c=n_c$.  The recovery-validity inequalities imply, on sampled chunks,
\[
 1\le s_c\le n_c,\qquad s_c\ge(S/n)n_c,\qquad s_c\ge\kappa.
\tag{8}\label{eq:sampling-properties}
\]

Every chunk caches $x_{n,c}$, the restriction of
$C_{n,c}=\mathsf{Enc}_{n_c}(x_{n,c})$ to $W_{\ell_c}$, the full word
$C_{n,c}$, and typed Merkle trees for all three arrays.  Its active data word is
\[
 D_{n,c}=\begin{cases}
 x_{n,c},&c\text{ tiny},\\
 C_{n,c}|_{W_{\ell_c}},&c\text{ mini},\\
 C_{n,c},&c\text{ sampled}.
 \end{cases}
\tag{9}\label{eq:active-word}
\]
Every tree, oracle, graph, and commitment namespace includes the full tag
$\iota_{n,c}$ from \Cref{eq:chunk-data}; this separates equal-size chunks in
the two forests.  Define the data-type tag
\[
 t^D_c=\begin{cases}
 \mathsf{raw},&c\text{ tiny},\\
 \mathsf{mini},&c\text{ mini},\\
 \mathsf{full},&c\text{ sampled},
 \end{cases}
 \qquad
 \gamma^D_c=(\mathsf{data},\iota_{n,c},t^D_c,m_c),
 \qquad
 \gamma^h_{c,\mathit{id}}
   =(\mathsf{label},\iota_{n,c},\mathit{id},4\tau_c).
\tag{T1}\label{eq:type-tags}
\]
Let $\mathit{ser}_{\rm pp}$ be the finite prefix-free serialization of the
additive-code tables, $\mathit{pp}_{\rm DRG}$, the commitment reference string,
the Merkle-hash identifier, the descriptions of the $G,H,FS$ domain
namespaces, and the scalar parameters.  It contains neither random-oracle
tables nor the identifier being defined.  Let $\mathit{id}_{\rm pp}$ be the
$\lambda_{\rm vc}$-bit hash, under a dedicated setup separator, of
$\mathit{ser}_{\rm pp}$.  The full setup remains ambient public data and is not
copied into each digest.  For
$0\le\ell\le L_{\max}$, define
\[
 \mathsf{slot}_{n,f,\ell}=
 \begin{cases}
  (\iota_{n,(f,\ell)},\varphi_{n,(f,\ell)}),&(f,\ell)\in\mathcal K(u),\\
  \bot,&(f,\ell)\notin\mathcal K(u),
 \end{cases}
\]
where $\varphi_{n,c}$ is the typed Merkle root of $D_{n,c}$.  The public digest
is the canonical serialization
\[
 \varphi_n=\left\langle
 (\mathsf{slot}_{n,f,\ell})_{f\in\{0,1\},\,0\le\ell\le L_{\max}},
 (n,n_0,N,S,B,\kappa,R,w,\Delta,\zeta_0,\mathit{id}_{\rm pp},
      (t^D_c)_{c\in\mathcal K(u)})\right\rangle.
\tag{T2}\label{eq:public-digest}
\]
Its parser recomputes $u=n-n_0$ and $\mathcal K(u)$, checks every active tag
and root and every inactive $\bot$ slot, and verifies the scalar header and
setup identifier.  If $z_0=|\langle\zeta_0\rangle|$, then
\[
 L_{\varphi_n}=O((1+\log N)(\lambda_{\rm vc}+\log N)+\log B+z_0).
\tag{T3}\label{eq:digest-length}
\]

The typed vector commitment has exact interfaces
\[
 \mathsf{VC.Digest}(\mathit{crs},\gamma,m,A)\to(\phi,\mathit{aux}),
 \qquad
 \mathsf{VC.Open}(\mathit{crs},\gamma,m,\mathit{aux},Q)\to\pi,
\]
\[
 \mathsf{VC.Vf}(\mathit{crs},\gamma,m,\phi,Q,A[Q],\pi)
 \to\{0,1\}.
\tag{T4}\label{eq:vc-interface}
\]
For every legal $\iota=(f,a,\ell,2^\ell)$ and $1\le\tau\le2N$, Setup exposes
independent persistent families
\[
 G_\iota:\{0,1\}^{*}\times\mathbb N\to[R2^\ell],\qquad
 H_{\iota,\tau}:\{0,1\}^{*}\to\mathbb F,
 \qquad
 FS_{\iota,\tau}:\{0,1\}^{*}\to[4\tau]^\kappa,
\tag{T5}\label{eq:oracle-domains}
\]
and the parent map
\[
 \operatorname{Par}_{\iota,\tau}(q)=
 \mathsf{DRG.Parents}(\mathit{pp}_{\rm DRG},1^\lambda,4\tau,\zeta_0,
      (\iota,\tau),q).
\tag{T6}\label{eq:parent-domain}
\]
Thus $G_c=G_{\iota_{n,c}}$, while $H_c$, $FS_c$, and
$\operatorname{Par}_c$ are indexed by the pair
$(\iota_{n,c},\tau_c)$.  This second component is mandatory because the active
sampling length of a persistent chunk may change.  For sampled $c$, put
$g_{c,\mathit{id}}(j)=G_c(\mathit{id},j)$ for $j\in[s_c]$; these values are
independent uniform coordinates sampled with replacement.  On a mini chunk
$g_{c,\mathit{id}}(j)=j$ for $j\in[2n_c]$.

Put $\mathcal Q_c=\{3\tau_c,\ldots,4\tau_c-1\}$.  The node stores
$d_{c,\mathit{id}}[j]=D_{n,c}[g_{c,\mathit{id}}(j)]$ and the label array
\[
 h_{c,\mathit{id}}[q]=
 \begin{cases}
 H_c((\varphi_n,\mathit{id}),q,
       h_{c,\mathit{id}}[\operatorname{Par}_c(q)]),&q<3\tau_c,\\
 H_c((\varphi_n,\mathit{id}),q,
       h_{c,\mathit{id}}[\operatorname{Par}_c(q)])
       \mathbin\oplus d_{c,\mathit{id}}[q-3\tau_c],&q\in\mathcal Q_c.
 \end{cases}
\tag{T7}\label{eq:labels}
\]
Here the parent tuple is ordered and has size at most $\Delta$.  Let
$(\chi_{c,\mathit{id}},\mathit{aux}^h_{c,\mathit{id}})$ be the result of
committing to this complete label array under $\gamma^h_{c,\mathit{id}}$.
For any finite tuple $Q=(q_1,\ldots,q_r)\in[4\tau_c]^r$, put
\[
 U_c(Q)=\bigcup_{j=1}^r
   (\{q_j\}\cup\operatorname{Par}_c(q_j)).
\tag{T8}\label{eq:neighborhood}
\]

The offline predicate is as follows.  Define
\[
 F_c=FS_c(\chi_{c,\mathit{id}},\mathit{id}),\qquad
 J_c=\{q-3\tau_c:q\in F_c\cap\mathcal Q_c\},\qquad
 K_c=\{g_{c,\mathit{id}}(j):j\in J_c\}.
\tag{T9}\label{eq:offline-sets}
\]
The prover supplies $h^0_c[U_c(F_c)]$, occurrence-indexed values
$d^0_c[J_c]$, and separate label and data proofs.  The verifier first checks
that occurrences mapped to one coordinate have equal values and forms the
coordinate-indexed array $\overline d^0_c[K_c]$.  It then verifies
\[
 \mathsf{VC.Vf}(\mathit{crs},\gamma^h_{c,\mathit{id}},4\tau_c,
 \chi_{c,\mathit{id}},U_c(F_c),h^0_c[U_c(F_c)],\pi^{0,h}_c)=1,
\]
\[
 \mathsf{VC.Vf}(\mathit{crs},\gamma^D_c,m_c,
 \varphi_{n,c},K_c,\overline d^0_c[K_c],\pi^{0,D}_c)=1,
\tag{T10}\label{eq:offline-openings}
\]
and checks \Cref{eq:labels} at every $q\in F_c$, using the supplied
occurrence value when $q\in\mathcal Q_c$.

After all offline predicates accept, the online verifier independently samples
$Q_c\in\mathcal Q_c^\kappa$, receives $h^e_c[U_c(Q_c)]$ and $\pi^e_c$, and
checks exactly
\[
 \mathsf{VC.Vf}(\mathit{crs},\gamma^h_{c,\mathit{id}},4\tau_c,
 \chi_{c,\mathit{id}},U_c(Q_c),h^e_c[U_c(Q_c)],\pi^e_c)=1.
\tag{T11}\label{eq:online-opening}
\]
A tiny chunk instead opens its complete raw word under $\gamma^D_c$.  Every
verification passes the tag, length, aligned index set, values, and proof
explicitly.  All chunks run in parallel, with at most $s_{\rm lb}(n)$
sequential label-oracle dependency rounds from receipt of the identity through
the online response.

For later rewinding, let $\Xi$ denote the typed public audit context specified
above together with the fixed persistent functions, and use the configuration
space $\Gamma_{\mathcal B}$ from \Cref{sec:preliminaries}.  Let
$\mathsf{Stat}=\{\mathsf{accept},\mathsf{reject},\mathsf{halt}\}$ and
$\mathsf{Id}=\{0,1\}^{*}$.  Let $\mathfrak O_{\rm off}$ and
$\mathfrak O_{\rm on}$ be the finite spaces of the typed records specified in
\Cref{eq:offline-sets,eq:offline-openings,eq:online-opening}.  For an active
chunk, set
$\mathcal C_c=\mathcal Q_c^\kappa$ when it is sampled or mini and
$\mathcal C_c=\{[n_c]\}$ when it is tiny.  With $\mathsf{Com}_c$ denoting the
label-commitment space, the transition interfaces are
\[
 \mathsf{NextId}_\Xi:\Gamma_{\mathcal B}\to
  (\mathsf{Id}\times\Gamma_{\mathcal B})
  \cup(\{\mathsf{halt}\}\times\Gamma_{\mathcal B}),
\]
\[
 \mathsf{Offline}_\Xi:\Gamma_{\mathcal B}\times\mathsf{Id}\to
  \mathsf{Stat}\times\Gamma_{\mathcal B}\times
  (\mathfrak O_{\rm off}\cup\{\bot\}),
\]
\[
 \mathsf{Online}_\Xi:\Gamma_{\mathcal B}\times\mathsf{Id}
 \times\prod_{c\text{ active}}\mathcal C_c
 \times\prod_{c\text{ sampled or mini}}\mathsf{Com}_c
 \to\mathsf{Stat}\times\Gamma_{\mathcal B}\times
 (\mathfrak O_{\rm on}\cup\{\bot\}).
\tag{T12}\label{eq:transition-signatures}
\]
The first map runs until the next identity submission or halt and returns the
exact configuration immediately after submission and before any offline
message.  The second executes all predicates in
\Cref{eq:offline-sets,eq:offline-openings}; it returns the verified
record, including every $\chi_{c,\mathit{id}}$, exactly on acceptance.  The
third sends the explicit challenges, executes \Cref{eq:online-opening} and
the tiny predicates, and returns the exact post-interaction configuration and
every received record on a nonhalting branch.  Rejection, halt, malformed, and
missing-record branches return their named status and $\bot$ record as
appropriate.  For fixed random functions all three maps are deterministic,
and each has cost at most $C_{\rm tr}$ from \Cref{eq:transition-cost}.

The publisher state is $(n_0,u,\mathcal F_0,\mathcal F_1)$.  Forest
$\mathcal F_0$ represents the fixed initial sequence with origin $0$, and
$\mathcal F_1$ represents the suffix with origin $n_0$.  Within a forest of
local length $v$, a binary-carry level $\ell$ has local start
$\sum_{j\in H(v),j>\ell}2^j$ and the corresponding full global interval tag.

\begin{algbox}
\textbf{\uline{\textsc{TwoForest}}}: The tagged
Setup--Init--Join--Update--Audit protocol suite.
\phantomsection\label{alg:twoforest}

\uline{Input}: Setup receives $1^\lambda$ and a publisher-admissible public
parameter tuple from \Cref{def:admissible};
Init receives $(S,\mathit{DB}_{n_0})$; Join receives a valid publisher state,
its digest, and $\mathit{id}$; Update receives a valid state, one block $b_n$,
and a participant set; Audit receives a digest, an identity, and a purported
node state.

\uline{Output}: Respectively public parameters; the initial publisher state
and digest; a current node state; the updated publisher state, digest, and
refreshed honest-node states, with the publisher state maintained in place; or
$\mathsf{accept}/\mathsf{reject}$.

\uline{Invariant}: Each forest is the tagged dyadic decomposition of its own
sequence, and every active chunk stores the three code representations and
typed trees defined above.

\begin{enumerate}[leftmargin=*,label=\arabic*.]
 \item \textsc{Setup}: invoke
 $\hyperref[alg:additive]{\textsc{AdditiveEC}}(\mathsf{setup},(B,N,R))$,
 initialize the typed Merkle
 commitment, and publish the independently tagged shard, label,
 Fiat--Shamir, and graph namespaces.
 \item \textsc{Init}: starting with empty $\mathcal F_0$, append
 $b_0,\ldots,b_{n_0-1}$.  For each append, create the three tagged level-zero
 objects, merge through the trailing occupied levels using
 $\hyperref[alg:additive]{\textsc{AdditiveEC}}
 (\mathsf{merge},\allowbreak(t,\allowbreak C^{(0)},\allowbreak C^{(1)}))$,
 and build the
 three tagged parent trees.  Freeze the resulting $\mathcal F_0$, set
 $\mathcal F_1=\varnothing$, select the active representations in
 \Cref{eq:active-word}, and serialize $\varphi_{n_0}$.
 \item \textsc{Join}: parse the digest and recompute all chunks and classes.
 Send authenticated data symbols for each sampled or mini shard and the full
 raw word for each tiny chunk.  Verify those openings, compute and commit the
 labels in \Cref{eq:labels}, and store the complete label arrays, required
 data, authentication paths, and offline records.
 \item \textsc{Update}: append $b_n$ only to $\mathcal F_1$.  Create its
 tagged singleton; while the old suffix length has a trailing occupied level,
 merge that older level with the newer carry, build the three parent trees,
 and release the children.  Leave $\mathcal F_0$ unchanged, increment $u$,
 select the cached current representations, serialize $\varphi_{n+1}$, and
 invoke Join independently for each displayed participant.
 \item \textsc{Audit}: verify the public header and offline records, sample
 the independent tuples $Q_c$, verify all openings on \Cref{eq:neighborhood}
 and every tiny raw opening, and accept exactly when all typed predicates and
 the common response-round bound hold.
\end{enumerate}
\end{algbox}

Correctness of the five entry points follows from the common forest invariant;
the same invariant also exposes the carry cost.

\begin{lemma}[Two-forest protocol correctness and publisher accounting]
\label{lem:publisher}
Under publisher admissibility, \hyperref[alg:twoforest]{\textsc{TwoForest}}
is honestly correct.  Init uses $O(BRn_0\log N)$ simplified work.  If the old
suffix length $u$ has $q$ trailing one-bits, one node-independent Update uses
\[
 O(BR2^q+B\log N+L_{\varphi_{n+1}})
\tag{13}\label{eq:one-update}
\]
simplified work.  For every $U\in[1,N-n_0]$, the first $U$ Updates use
\[
 O(BRU\log N+Uz_0)
\tag{14}\label{eq:aggregate-update}
\]
simplified work, or
$O((M(B)+B)RU\log N+Uz_0)$ bit work.  Publisher storage is $O(BRN)$, a full
tree has depth $O(\log N+\log R)$, and honest node state is $O(BS)$ under the
block-size hypothesis in \Cref{thm:main}.  The bounds include each serialized
digest and exclude Setup, Init, and identity-dependent Join output from the
post-Init maintenance measure.  More precisely, one honest Join at length $n$
materializes
\[
 L_{\rm J}(n,\mathit{id})=
 O(BS+\lambda_{\rm vc}S(\log N+\log R)+L_{\varphi_n})
\tag{U1}\label{eq:join-output}
\]
bits of transcript and node state.  Thus $d$ participants contribute
$O(dBS)$ bits under that block-size hypothesis, and the single public digest
makes the complete externally materialized output $O((d+1)BS)$.  The mutable
publisher hierarchy is maintained in place and is not reserialized in this
output quantity.
\end{lemma}

\begin{proof}
Consider one forest with origin $o$.  A singleton starts at $o+v$.  If the old
local length has $q$ trailing ones, the level-$j$ message is older than and
adjacent to the current carry for every $j<q$.  The rule in \Cref{eq:merge}
therefore produces their chronological concatenation.  The parent begins at
the old child's global start, so induction gives exactly the interval tag in
\Cref{eq:chunk-data}.  Higher levels remain unchanged.  This proves the
invariant.  It also proves the per-call coding bound because the geometric
sum of created word and tree sizes is $O(R2^q)$.

For compatibility with the new-before-old merge convention in the comparison
iDDA construction~\cite{ShiSilverMu2026}, no cached word must be changed.  If
$\mathsf{Rev}_m$ reverses a word and
$\widetilde{\mathsf{Enc}}_m(y)=\mathsf{Enc}_m(\mathsf{Rev}_m(y))$, then
\[
 \mathsf{Rev}_{2m}(u\Vert v)
   =\mathsf{Rev}_m(v)\Vert\mathsf{Rev}_m(u).
\]
Thus a protocol-facing new-before-old merge is conjugate to the chronological
old-before-new merge used here.  The same reversal applies to raw tiny chunks.

The initial forest is built once and is never touched by Update.  During $U$
suffix appends, a suffix chunk of size $2^\ell$ is created
$\lfloor U/2^\ell\rfloor$ times.  Hence
\[
 \sum_{\ell=0}^{\lfloor\log_2U\rfloor}
 O(BR2^\ell)\left\lfloor\frac U{2^\ell}\right\rfloor
 =O(BRU(1+\log U))=O(BRU\log N).
\]
By \Cref{eq:digest-length}, $R\widehat N\le2^B$, and
$\lambda_{\rm vc}\le B$, serializing all $U$ digests costs
$O(BU\log N+Uz_0)$, which proves \Cref{eq:aggregate-update}.  Replacing the
field-operation factor $B$ by $M(B)+B$ gives the bit-work bound.  The same
level count applied to the $n_0$ Init appends gives $O(BRn_0\log N)$.

The active full words contain
$R\sum_{c\in\mathcal K(u)}n_c=Rn\le RN$ field elements.  The raw words,
redundancy-two restrictions, typed trees, and a constant number of carry
buffers change this by only a constant factor.  This proves storage, and the
largest tree has at most $RN$ leaves.

For honest-node space, \Cref{eq:sample-size} and \Cref{lem:partition} give
\[
 \sum_{c\text{ sampled}}s_c
 \le S+O(|\mathcal K(u)|(s_{\rm lb}(n)+1))=O(S),
\]
using \Cref{eq:recovery-validity}, $n\ge n_0\ge\lambda$, and
$N\le\lambda^{O(1)}$.  In each forest the total size of mini chunks is less
than $2s_{\rm lb}(n)$ and the total size of tiny chunks is less than
$2\kappa$, so the two forests still contribute $O(S)$ block values.  Merkle
paths use $O(\lambda_{\rm vc}S(\log N+\log R))=O(BS)$ bits.  The label
recurrence and typed commitment correctness make every honest offline and
online opening verify.  Finally, all three representations already exist, so
a class change only selects another cached root; its participant-specific
refresh is precisely the separately charged Join call.  This completes the
correctness and resource proof.  The same count includes every value and
authentication path written by Join and the $L_{\varphi_n}$-bit public header,
which proves \Cref{eq:join-output}; summing it over $d$ independent sessions
gives the stated output bounds.
\end{proof}

The remainder of the proof establishes the paired recovery clause.

\subsection{Exact challenges and accepting-spine extraction}

Fix a recovery-valid length $n=n_0+u$, meaning that the inequalities in
\Cref{eq:recovery-validity} hold.  For a sampled chunk define the public
allocation
\[
 r_c=\min\left\{\mu,
    \left\lceil\frac{n_c}{\vartheta s_c}\right\rceil\right\},\qquad
 c_c=\min\{n_c,\lceil\vartheta\mu s_c\rceil\},\qquad
 d_c=n_c-c_c.
\tag{15}\label{eq:allocation}
\]
For mini and tiny chunks put $(r_c,c_c,d_c)=(1,n_c,0)$.  This allocation is a
function only of public data.  Let
$r_{\max}=\max(\{1\}\cup\{r_c:c\text{ sampled}\})$.

Set $T_{\rm ext}=\lceil\lambda^{\log_2\lambda}\rceil$.  With
$L=\lfloor\log_2\lambda\rfloor$, define
\[
 a_\epsilon=\lfloor\eta\kappa L\rfloor,
\quad
 e_\epsilon=\max\left\{1,
 \left\lfloor\frac14\min\{\lfloor\sqrt{a_\epsilon}\rfloor,L^2\}
 \right\rfloor\right\},
\quad K_{\rm sp}=2^{e_\epsilon}.
\tag{16}\label{eq:spine-count}
\]
For each sampled or mini chunk let
$b_c=\max\{1,\lceil\log_2\tau_c\rceil\}$.  There are two disjoint public tag
sets.  The spine set contains
$(\mathsf{spine},a,i,c,v)$ for $a\in[K_{\rm sp}]$, $i\in[\mu]$, and
$v\in[\kappa]$, with $c$ ranging over all sampled or mini chunks.  The raw set contains trial tags
$(\mathsf{trial},i,t,c,v)$ for $i\in[r_{\max}]$,
$t\in\{1,\ldots,T_{\rm ext}\}$, and canonical tags
$(\mathsf{canonical},i,c,v)$, again with $c$ sampled or mini and
$v\in[\kappa]$.  Tiny-chunk
challenges are deterministic full openings.

For either finite tag set $\mathcal I_X$, put
\[
 A_X=\lambda+\lceil\log_2\max\{1,|\mathcal I_X|\}\rceil+2.
\]
For a tag belonging to chunk $c$, allocate $A_X$ disjoint $b_c$-bit blocks.
Interpret a block as an integer in $[2^{b_c}]$ and select the first block below
$\tau_c$; fail if none exists.  Adding $3\tau_c$ converts the selected integer
to an element of $\mathcal Q_c$.  With $c(\alpha)$ denoting the chunk in tag
$\alpha$, the exact substring lengths are
\[
 D_X=A_X\sum_{\alpha\in\mathcal I_X}b_{c(\alpha)},\qquad
 \rho_{\rm sp}\in\{0,1\}^{D_{\rm sp}},\qquad
 \rho_{\rm raw}\in\{0,1\}^{D_{\rm raw}},
\tag{S1}\label{eq:coin-lengths}
\]
and the shared string is
$\rho=\rho_{\rm sp}\Vert\rho_{\rm raw}$ of length
$D_{\rm sp}+D_{\rm raw}$.

\begin{algbox}
\textbf{\uline{\textsc{TaggedChallengeSample}}}: Parse finite shared coins
into exact independent audit challenges.
\phantomsection\label{alg:challenge-sample}

\uline{Input}: A public sampling state $\mathfrak X$ containing the tagged
chunks, $(\tau_c)_c$, $\kappa$, $\mu$, $r_{\max}$, $K_{\rm sp}$, and
$T_{\rm ext}$, and a bit string $\rho=\rho_{\rm sp}\Vert\rho_{\rm raw}$ of
the exact length in \Cref{eq:coin-lengths}.

\uline{Output}: The complete spine, raw-trial, and canonical challenge tables,
or $\mathsf{fail}_{\rm sp}$ or $\mathsf{fail}_{\rm raw}$.

\uline{Guarantee}: Conditional on neither failure, every selected coordinate
is exactly uniform in its chunk range, all coordinates are mutually
independent, and they are independent of the receiver experiment.

\begin{enumerate}[leftmargin=*,label=\arabic*.]
 \item In lexicographic tag order, partition $\rho_{\rm sp}$ and
 $\rho_{\rm raw}$ into the prescribed disjoint blocks; reject a string of any
 other length.
 \item For each spine tag, set a spine-failure flag if none of its
 $A_{\rm sp}$ blocks is below $\tau_c$; otherwise record the first such value.
 \item Independently apply the same rule with $A_{\rm raw}$ to every raw and
 canonical tag and set a raw-failure flag when necessary.  Thus both flags are
 defined on every coin-string outcome.
 \item If a flag is set, return its named failure value, giving spine failure
 priority only for the output syntax.  Otherwise add $3\tau_c$ to each
 recorded value and output the named tables.
\end{enumerate}
\end{algbox}

The bounded-rejection construction has an exact distribution, as recorded
next.

\begin{lemma}[Exact bounded-rejection challenge sampling]
\label{lem:challenge-sampling}
For uniform correctly sized $\rho$, each of the two failure events in
\hyperref[alg:challenge-sample]{\textsc{TaggedChallengeSample}} has probability
at most $2^{-\lambda-2}$.  Conditional on success, its distribution has the
stated exact product form.  Its running time is linear in its coin-string
length, and that length is quasi-polynomial under the parameter conditions of
\Cref{thm:main}.
\end{lemma}

\begin{proof}
Because $2^{b_c}<2\tau_c$ unless $\tau_c=1$, a single block succeeds with
probability at least $1/2$.  One tag therefore fails with probability at most
$2^{-A_X}$, and a union bound over $|\mathcal I_X|$ tags is at most
$2^{-\lambda-2}$.  Conditional on success for one tag, its first accepted
block is uniform on $[\tau_c]$.  Different tags use disjoint bits, and their
success events also depend on disjoint bits, so conditioning on global success
preserves independence.  The coin string is independent of the experiment.
Finally, $\tau_c\le2N$, while the two tag sets have quasi-polynomial size once
$K_{\rm sp}$ and $T_{\rm ext}$ do; scanning every allocated block is linear in
their total length.
\end{proof}

We next state precisely the local graph-extraction property used by the rewind
procedure.  It isolates the verified-occurrence conclusion supplied by the
iDDA extraction analysis~\cite{ShiSilverMu2026} for the exact typed Audit
relation required here.  The chosen graph and commitment instantiation must
satisfy this explicit hypothesis; the transcript-free coupling and collision
arguments that follow are proved locally.

\begin{definition}[Local verified-extraction contract]
\label{def:local-extraction}
Fix the typed audit relation in \Cref{eq:labels,eq:neighborhood}.
Starting at a saved configuration immediately after an identity is submitted
and before its offline message, a raw trial restores that machine
configuration, retains all persistent random functions, executes the offline
transition, and, if it accepts, executes the online transition with an
independent challenge.  Only label values in an accepted online typed opening,
and raw values in an accepted tiny-chunk online opening, are written.  Offline
Fiat--Shamir openings are checked but never enter this aggregate.  For each
coordinate, later accepted online writes replace earlier writes.

More explicitly, let $b^{\rm off}_{i,t}$ and $b^{\rm on}_{i,t}$ be the two
trial statuses, let $Q^e_{c,i,t}$ be its named online challenge, and let
$h^e_{c,i,t}$ be the received online label array.  Define
\[
 \mathsf W_{c,i,t}[q]=
 \begin{cases}
 h^e_{c,i,t}[q],&b^{\rm off}_{i,t}=b^{\rm on}_{i,t}=\mathsf{accept}
       \text{ and }q\in U_c(Q^e_{c,i,t}),\\
 \bot,&\text{otherwise},
 \end{cases}
\]
and, coordinatewise,
\[
 \mathsf H^{(0)}_{c,i}[q]=\bot,\qquad
 \mathsf H^{(t)}_{c,i}[q]=
 \begin{cases}
  \mathsf W_{c,i,t}[q],&\mathsf W_{c,i,t}[q]\ne\bot,\\
  \mathsf H^{(t-1)}_{c,i}[q],&\mathsf W_{c,i,t}[q]=\bot.
 \end{cases}
\tag{R1}\label{eq:verified-write-recurrence}
\]
The final verified label array is
$\widetilde h_{c,i}=\mathsf H^{(T_{\rm ext})}_{c,i}$.  On a tiny chunk, define
$\mathsf W^{\rm tiny}_{c,t}[q]$ to be the received raw value exactly when both
trial statuses accept, and $\bot$ otherwise; applying the same recurrence
gives its final partial raw array.  Thus every offline rejection, online
rejection, halt, or missing record makes no write.

For a sampled chunk, selected identity $\mathit{id}_i$, and occurrence
$j\in[s_c]$, put $q=3s_c+j$.  If the final
verified label array contains $q$ and all its parents, its candidate is
\[
 \left(g_{c,\mathit{id}_i}(j),
 \widetilde h_{c,i}[q]\mathbin\oplus
 H_c((\varphi_n,\mathit{id}_i),q,
   \widetilde h_{c,i}[\operatorname{Par}_c(q)])\right);
\tag{17}\label{eq:candidate}
\]
otherwise it is $\bot$.  A candidate is \emph{usable} if its value equals the
committed codeword at its displayed coordinate.  The procedure never tests
usability.  The mini-chunk definition is identical with the injective map
$j\mapsto j$ on $[2n_c]$, and a tiny chunk records the latest accepted full
raw opening.

The local verified-extraction contract asserts the following, with its
internal extraction slack instantiated as $2\eta$.  On the joint
probability space of the receiver experiment, persistent random functions, and
shared coins, there is an event $\mathsf{Good}_{\rm loc}$ contained in the
event that both coin parsers succeed and an accepting spine is found, such that
\[
 \Pr[\mathsf{Pass}\wedge\neg\mathsf{Good}_{\rm loc}]
 \le 2^{-\lambda-1}+\frac1{K_{\rm sp}+1}
 +K_{\rm sp}q_{\rm loc}
 \left(\frac{8N}{T_{\rm ext}+1}+(1-\eta)^\kappa\right)
 +\frac{K_{\rm sp}\mu}{T_{\rm ext}+1}
 +\nu_{\rm bind}(\lambda),
\tag{18}\label{eq:local-tail}
\]
and, on $\mathsf{Good}_{\rm loc}$, every sampled chunk and every selected
identity $i<r_c$ has at least
$(1-4\eta)s_c$ usable occurrences; each mini chunk has at least
$(1-4\eta)2n_c$ usable, distinct-coordinate occurrences for identity zero;
and each tiny chunk is recovered completely and correctly.  Here
\[
 q_{\rm loc}=12\mu(1+\lfloor\log_2N\rfloor)
 (C_{\rm tr}+N+\kappa+1),
\]
and $\nu_{\rm bind}$ is the advantage of one quasi-polynomial commitment
binding adversary.  This assertion is required uniformly for the exact
transitions in \Cref{eq:offline-openings,eq:online-opening}, the
sampled word $\mathsf{Enc}_{n_c}(x_{n,c})$, its $W_{\ell_c}$ mini restriction,
the full-opening tiny word $x_{n,c}$, and every full instance and length tag
$(\iota_{n,c},\tau_c)$.  Together with the typed commitment and parent
interfaces, these are the complete compatibility premises of the contract.
\end{definition}

An accepting spine is a list
$\mathcal S=((\mathit{id}_i,\sigma_i))_{i\in[\mu]}$ obtained on one execution
in which all $\mu$ audits accept, the identities are distinct, and
$\sigma_i$ is the complete configuration immediately after submission of
$\mathit{id}_i$ and before its offline message.  From $\mathcal S$, define a
checkpoint wrapper that returns these saved pairs in order, delegates each
offline or online transition to the original continuation from the saved
configuration, and advances exactly once after every local audit branch.  Its
next saved checkpoint is therefore independent of audit history supplied to
the wrapper.

\begin{algbox}
\textbf{\uline{\textsc{AdaptiveRawExtract}}}: Find an accepting adaptive path
and extract verified candidate symbols at its saved checkpoints.
\phantomsection\label{alg:adaptive-raw}

\uline{Input}: A public extraction state $\mathfrak X$ containing the typed
audit context, deterministic continuation $\mathcal B$, allocation $(r_c)_c$,
$K_{\rm sp}$, and $T_{\rm ext}$, and the shared coin string $\rho$.

\uline{Output}: Candidate pairs or $\bot$ for each sampled
$(c,i,j)$ with $i<r_c,j<s_c$, one length-$2n_c$ candidate array per mini
chunk, and one partial raw array per tiny chunk.

\uline{Invariant}: Machine configurations may be restored, but the same fixed
random functions and their lazy tables persist across every search and rewind.

\begin{enumerate}[leftmargin=*,label=\arabic*.]
 \item Invoke
 $\hyperref[alg:challenge-sample]{\textsc{TaggedChallengeSample}}
 (\mathfrak X,\rho)$.  On either
 sampler failure, return the correctly shaped all-$\bot$ record.
 \item For $a=0,\ldots,K_{\rm sp}-1$, restore the initial configuration of
 $\mathcal B$ and execute $\mu$ sequential audits with spine table $a$.
 Reject that path on a halt, a repeated identity, or any failed offline or
 online transition.  Save every pre-offline pair
 $(\mathit{id}_{a,i},\sigma_{a,i})$ on a fully accepting path and select the
 least accepting $a$.  Return the all-$\bot$ record if none exists.
 \item Form the checkpoint wrapper from the selected spine.  At its $i$th
 saved checkpoint, perform the $T_{\rm ext}$ raw trials in increasing trial
 order, always restoring $\sigma_i$ and retaining the oracle tables.  If an
 offline transition rejects or halts, do not invoke its online transition and
 make no write; if an online transition rejects or halts, likewise make no
 write.  Aggregate only accepted online verified label and tiny-word writes by
 \Cref{eq:verified-write-recurrence}.  Mini and tiny chunks use only
 checkpoint $i=0$; later checkpoints make no assignment to those outputs.
 \item Form sampled and mini candidates by \Cref{eq:candidate}.  After the
 trials at checkpoint $i$, use its separately tagged canonical challenge in
 one delegated audit so that the wrapper advances to checkpoint $i+1$.  A
 canonical offline rejection or halt advances immediately; after a canonical
 offline acceptance, the online branch advances regardless of its terminal
 status.
 Return all candidate and tiny arrays after $r_{\max}$ checkpoints.
\end{enumerate}
\end{algbox}

We now justify the accepting-path coupling, shared-output property, and total
cost of this extraction procedure.

\begin{lemma}[Transcript-free accepting-spine extraction]
\label{lem:spine}
Under the recovery conditions of \Cref{thm:main}, the following hold.
\begin{enumerate}[label=(\roman*),leftmargin=*]
 \item Jointly with $\mathsf{Pass}$, the probability that no accepting spine
 is found is at most
 $2^{-\lambda-1}+1/(K_{\rm sp}+1)$.
 \item Two callers with the same public context, continuation state,
 persistent functions, and $\rho$ select the same spine and obtain the same
 output from \hyperref[alg:adaptive-raw]{\textsc{AdaptiveRawExtract}}.  This
 computation uses no experiment transcript once the continuation state is
 fixed.
 \item The failure probability of the usable-occurrence conclusions in
 \Cref{def:local-extraction}, together with the two sampler failures and the
 spine miss, is negligible.
 \item The algorithm runs in time quasi-polynomial in
 $(\lambda,t_{\mathcal A})$.
\end{enumerate}
\end{lemma}

\begin{proof}
Condition on the public context, the initial continuation configuration, and
the complete persistent random functions, and then on successful parsing of
both independent coin substrings.  Cache membership is hidden, so a full path
is a deterministic function of its fresh challenge table.  If one such table
produces $\mu$ accepting distinct-identity audits with probability $p$, then
the received experiment table and the $K_{\rm sp}$ search tables are
independent samples from the same conditional distribution.  Therefore
\[
 \Pr[\mathsf{Pass}\wedge\text{no searched path accepts}]
 =p(1-p)^{K_{\rm sp}}\le\frac1{K_{\rm sp}+1}.
\]
Averaging and adding both sampler failures proves (i).

All transitions and parsers are deterministic after fixing the displayed
data.  Both callers consequently choose the same least accepting path.  The
checkpoint wrapper invokes the original continuation at precisely the saved
configurations and makes those checkpoints audit-history independent.  The
last-write recurrence and candidate map are deterministic, proving (ii).

It remains to check the variable-slack balance.  Put
$y=\eta\kappa/L$.  Since
$\overline\epsilon>\epsilon/2$ and
$\epsilon\kappa=\omega(\log\lambda)$, one has $y\to\infty$.  From
\Cref{eq:spine-count},
\[
 e_\epsilon=\frac L4\min\{\sqrt y,L\}+O(1)=\omega(L),
 \qquad e_\epsilon=o(\eta\kappa),
 \qquad e_\epsilon\le L^2/4+1.
\tag{19}\label{eq:spine-balance}
\]
Hence $K_{\rm sp}^{-1}$ is negligible,
$\log K_{\rm sp}=o(\eta\kappa)$, and
$K_{\rm sp}\mu/(T_{\rm ext}+1)$ is negligible.  Moreover
$(1-\eta)^\kappa\le\exp(-\eta\kappa)$, while
$q_{\rm loc}$ is polynomial.  Thus every term in \Cref{eq:local-tail} is
negligible; $\nu_{\rm bind}$ is negligible by quasi-polynomial binding.  With
\Cref{lem:challenge-sampling}, this proves (iii).

Finally, $K_{\rm sp},T_{\rm ext}\le2^{O((\log\lambda)^2)}$, and all transition
and record sizes are polynomial in the displayed parameters.  The number of
transitions is $O(K_{\rm sp}\mu+\mu T_{\rm ext})$, so the claimed running time
follows from \Cref{eq:transition-cost}.
\end{proof}

The selected identities depend on oracle answers.  We therefore need an
expansion estimate that remains valid after deletions and an exposure argument
that applies it without conditioning on spine success.

\subsection{Collision expansion and completion}

\begin{lemma}[Deletion-resilient adaptive shard expansion]
\label{lem:collision}
Let $0<\eta\le1/40$, $\vartheta=1-13\eta$, $1\le s\le m$, and
$R\ge12\mathrm e/\eta$.  Let an adaptive process make at most $T$ distinct
queries to a random function
$G:\{0,1\}^{*}\to[Rm]^s$ whose fresh answer is an independent uniform
$s$-tuple.  After exposure, retain an arbitrary set
$A_x\subseteq[s]$ of size at least $(1-4\eta)s$ for each queried input $x$.
Put $r_* =\lceil m/(\vartheta s)\rceil$.  Except with probability at most
\[
 \sum_{r=1}^{\min\{T,r_*\}}\binom Tr4^{-\eta sr},
\tag{20}\label{eq:collision-failure}
\]
every $J$ of queried inputs satisfies
\[
 \left|\bigcup_{x\in J}\{G(x)[j]:j\in A_x\}\right|
 \ge\min\{m,(1-5\eta)s|J|\}.
\tag{21}\label{eq:collision-expansion}
\]
If $\eta s=\omega(\log\lambda)$ and
$\log_2\max\{2,T\}=o(\eta s)$, the failure probability is negligible.
\end{lemma}

\begin{proof}
Pad an early-stopping process with fresh dummy inputs.  Fix $r\le r_*$ query
positions and flatten their answers into $q=sr$ uniform draws from $[Rm]$.
Since $\vartheta>1/2$ and $s\le m$, one has $q<3m$.  Conditional on all prior
draws, the probability that the next draw repeats is less than $3/R$.  If at
least $\eta q$ draws repeat, some $d=\lceil\eta q\rceil$ positions witness this.
Iterated conditioning and a union bound give
\[
 \Pr[\text{at least }\eta q\text{ repeats}]
 \le\binom qd(3/R)^d
 \le\left(\frac{3\mathrm e}{R\eta}\right)^d
 \le4^{-\eta q}.
\]
Union-bounding over the $\binom Tr$ choices and all $r\le r_*$ proves
\Cref{eq:collision-failure}.  Deleting occurrences cannot increase their
inherited repeat count.  Thus at least $(1-4\eta)q$ retained occurrences lose
fewer than $\eta q$ new coordinates and occupy at least
$(1-5\eta)q$ coordinates.  For $|J|>r_*$, choose an $r_*$-subset and use
\[
 (1-5\eta)sr_*\ge
 \frac{1-5\eta}{1-13\eta}m\ge m.
\]
This proves \Cref{eq:collision-expansion}.  Finally, with
$z=T4^{-\eta s}$, the sum in \Cref{eq:collision-failure} is at most
$z/(1-z)$; its logarithm is $-(2-o(1))\eta s=-\omega(\log\lambda)$.
\end{proof}

The collision lemma must be invoked on a total exposure process, rather than
after selecting the potentially shard-oracle-dependent spine.

\begin{lemma}[Unconditional expansion for the selected spine]
\label{lem:deferred}
For every sampled chunk $c$, there is a padded adaptive exposure process whose
query set $\mathcal X_c^{\rm pad}$ contains exactly $p_{\rm sel}$ distinct
inputs to the vector-valued shard random function.  Write
$\mathsf{SampleOK}$ for success of both coin parsers and $\mathsf{SpineOK}$
for success of the accepting-spine search.  On
$\mathsf{SampleOK}\cap\mathsf{SpineOK}$, the set contains every selected
identity $\mathit{id}_0,\ldots,\mathit{id}_{r_c-1}$, and it satisfies
\[
 \log_2p_{\rm sel}=o(\eta s_c).
\tag{22}\label{eq:exposure-budget}
\]
Let $\mathsf{ExpBad}_c$ be the event that there are
$J\subseteq\mathcal X_c^{\rm pad}$ with
$1\le|J|\le\lceil n_c/(\vartheta s_c)\rceil$ and sets
$A_x\subseteq[s_c]$, $|A_x|\ge(1-4\eta)s_c$, such that
\[
 \left|\bigcup_{x\in J}\{G_c(x,j):j\in A_x\}\right|
 <(1-5\eta)s_c|J|.
\]
Without conditioning on sampler or spine success, this event has
probability at most
\[
 \nu_{{\rm exp},c}=
 \sum_{r=1}^{\min\{p_{\rm sel},\lceil n_c/(\vartheta s_c)\rceil\}}
 \binom{p_{\rm sel}}r4^{-\eta s_cr}
 =\operatorname{negl}(\lambda).
\tag{23}\label{eq:exposure-failure}
\]
Consequently, on
$\mathsf{SampleOK}\cap\mathsf{SpineOK}\cap
\mathsf{ExpBad}_c^{\mathrm c}$, any sets of usable occurrences
$A_i\subseteq[s_c]$ with $|A_i|\ge(1-4\eta)s_c$ satisfy
\[
 \left|\bigcup_{i<r_c}
 \{g_{c,\mathit{id}_i}(j):j\in A_i\}\right|
 \ge(1-5\eta)r_cs_c.
\tag{24}\label{eq:selected-expansion}
\]
\end{lemma}

\begin{proof}
The exposure reduction has four deterministic steps once its random function
is fixed.
\begin{enumerate}[leftmargin=*,label=\arabic*.]
 \item Implement the shard oracle by sampling its entire $s_c$-coordinate
 vector when a first argument is queried for the first time, while revealing
 to the machine only the requested coordinates.
 \item Expose every distinct first argument queried by the received experiment
 and its replay, all reached spine searches, and a pruned canonical checkpoint
 wrapper.  On $\mathsf{SampleOK}\cap\mathsf{SpineOK}$, also expose the selected
 $r_c$ identities.
 \item Let $q_{\rm pre}=q_{\rm pre}(\lambda)$ be a deterministic polynomial
 upper bound on the number of distinct shard-oracle first arguments made by
 the polynomial-time received experiment and its initialization replay.  Take
 $p_{\rm sel}$ to be a deterministic integer upper bound on
\[
 q_{\rm pre}+3(K_{\rm sp}+1)\mu C_{\rm tr}+\mu.
\]
 Such a bound exists independently of the experiment outcome.  Pad every
 outcome with the length-lexicographically first fresh dummy inputs and denote
 the resulting $p_{\rm sel}$-element input set by
 $\mathcal X_c^{\rm pad}$.
 \item Defer every shard-oracle query made only inside a restored local trial
 until after the preceding exposure.  Then execute each omitted trial from its
 saved configuration.
\end{enumerate}
The final reordering does not change any answer because the oracle is a fixed
function; cache membership is hidden, and discarded trial configurations are
not retained.  Thus it preserves the joint law of the experiment,
the search, the selected identities, and the raw output.  By
\Cref{eq:spine-balance}, polynomiality of the remaining factors, and
$s_c\ge\kappa$,
\[
 \log_2p_{\rm sel}\le\log_2K_{\rm sp}+O(\log\lambda)
 =o(\eta\kappa)=o(\eta s_c).
\]
Also $\eta s_c=\omega(\log\lambda)$.  Apply \Cref{lem:collision} unconditionally
with $(m,s,T)=(n_c,s_c,p_{\rm sel})$.  This proves
\Cref{eq:exposure-failure}; the simultaneous retained-set conclusion gives
\Cref{eq:selected-expansion} even though the sets $A_i$ are determined after
the rewinds.
\end{proof}

We now define the partial word that both recovery callers compute.  For every
integer coordinate $q\in[Rn_c]$, take the lexicographically first populated raw
candidate $(i,j)$ carrying coordinate $q$, and write its value at the evaluation
point $z_{\ell_c,q}$.  Write $Y_c$ for the resulting partial word,
$P_c=\operatorname{supp}(Y_c)$, and, only for analysis, let $e_c$ be the number
of supported values that disagree with $C_{n,c}$.  If $d_c>0$, let $I_c$ be
the first $d_c$ points of $V_{\ell_c}\setminus P_c$ in public order; if
$d_c=0$, let $I_c=\varnothing$.  For $v\in\mathbb F^{d_c}$,
$\mathsf{Fill}_c(Y_c,I_c,v)$ agrees with $Y_c$ on $P_c$, with $v$ on $I_c$,
and is $\bot$ elsewhere.  Mini candidates occupy their fixed first $2n_c$
evaluation points, and tiny candidates form a partial raw message.  None of
these operational maps computes $e_c$.

\begin{lemma}[Effective first-write support]
\label{lem:effective-support}
On
$\mathsf{Good}_{\rm loc}\cap
 \bigcap_{c\text{ sampled}}\mathsf{ExpBad}_c^{\mathrm c}$,
every sampled chunk satisfies
\[
 |P_c|-2e_c\ge
 c_c=\min\{n_c,\lceil\vartheta\mu s_c\rceil\}.
\tag{25}\label{eq:effective-support}
\]
The public complement $I_c$ exists on every raw record, not merely on this good
event.  Both paired recovery callers compute identical $Y_c,P_c,I_c$.
Furthermore every mini partial word has support size $m_c'$ and error count
$e_c'$ satisfying $m_c'-2e_c'\ge n_c$, and every tiny message is correct.
\end{lemma}

\begin{proof}
By definition, $\mathsf{Good}_{\rm loc}$ implies
$\mathsf{SampleOK}\cap\mathsf{SpineOK}$, so \Cref{lem:deferred} applies.
Fix a sampled chunk and write $q=r_cs_c$.  For each $i<r_c$, let $A_i$ be its
usable occurrences.  There are at most $4\eta q$ nonusable occurrences in
total.  By \Cref{eq:selected-expansion}, the usable occurrences occupy at
least $(1-5\eta)q$ distinct coordinates.  A nonusable occurrence can precede a
usable one and corrupt the selected value at at most one coordinate, and every
erroneous selected value is charged to a nonusable occurrence.  Hence
$e_c\le4\eta q$ and
\[
 |P_c|-2e_c\ge(1-5\eta)q-8\eta q
 =(1-13\eta)q=\vartheta r_cs_c.
\]
The left side is integral.  If $r_c=\mu$, this gives
$\lceil\vartheta\mu s_c\rceil$, capped by $n_c$.  If $r_c<\mu$, the definition
of $r_c$ gives $\vartheta r_cs_c\ge n_c$.  This proves
\Cref{eq:effective-support}.

If $d_c>0$, then $r_c=\mu$ and
$\mu s_c<n_c/\vartheta<2n_c$.  Thus $|P_c|<2n_c$.  Since $R\ge3$, at least
$Rn_c-|P_c|\ge n_c\ge d_c$ public evaluation points remain, so $I_c$ exists.
If $d_c=0$ the claim is immediate.  Determinism of the shared spine, raw
record, and public first-coordinate rule proves paired equality.

For a mini chunk, at least $(1-4\eta)2n_c$ of the fixed coordinate occurrences
are usable.  If $w_c'$ positions are missing, then
$e_c'+w_c'\le4\eta(2n_c)$, whence
\[
 m_c'-2e_c'=2n_c-w_c'-2e_c'
 \ge2n_c-2(e_c'+w_c')
 \ge(1-8\eta)2n_c\ge n_c.
\]
The tiny conclusion is part of the local extraction contract.
\end{proof}

The final algorithm sends only values at the public complement; in particular,
it never transmits support indices.

\subsection{Paired recovery and conclusion}

For a public allocation, serialize the complement values in the fixed order of
$\mathcal K(u)$, using exactly $B$ bits per field element.  The inverse parser
expects precisely $B\sum_{c\text{ sampled}}d_c$ bits and returns $\bot$ on any
other length.  The compressor may replay $\mathcal A_1$ from $\mathit{tr}$ to
recover $(n,\varphi_n,\mathit{DB}^{*},\mathit{st}_{\mathcal A})$; after this
replay, both callers form the same extraction state $\mathfrak X$.

\begin{algbox}
\textbf{\uline{\textsc{TwoForestRecovery}}}: Paired no-index compression and
error-and-erasure extraction.
\phantomsection\label{alg:recovery}

\uline{Input}: The compressor entry receives
$(1^\lambda,S,\mu,\mathit{tr};\rho)$ and is parameterized by the fixed
adversary $\mathcal A=(\mathcal A_1,\mathcal A_2)$.  The extractor entry
receives the continuation $\mathcal A_2(\mathit{st}_{\mathcal A})$ and
$(1^\lambda,n,S,\mu,\varphi_n,\mathit{DB}_{\rm short};\rho)$.  Both have the
same ambient public parameters and persistent random functions.

\uline{Output}: The compressor returns $\mathit{DB}_{\rm short}$; the extractor
returns $\mathit{DB}_{\rm ext}$ or $\bot$.

\uline{Guarantee}: On paired inputs, the two entries use the same raw record and
public complement sets.  The extractor has no transcript or original-database
input.

\begin{enumerate}[leftmargin=*,label=\arabic*.]
 \item The compressor replays $\mathcal A_1$ from $\mathit{tr}$ to recover
 $(n,\varphi_n,\mathit{DB}^{*},\mathit{st}_{\mathcal A})$; the extractor uses
 its explicitly supplied $n,\varphi_n$ and continuation.  Both entries then
 parse $\varphi_n$, compute $u=n-n_0$, $\mathcal K(u)$, the classes, and the
 allocation \Cref{eq:allocation}.  On a failed replay, malformed public data,
 or incorrectly sized $\rho$, the compressor returns the empty string and the
 extractor returns $\bot$.  Otherwise each invokes
 $\hyperref[alg:adaptive-raw]{\textsc{AdaptiveRawExtract}}(\mathfrak X,\rho)$
 and constructs
 $(Y_c,P_c,I_c)_c$ by the public first-coordinate rule.
 \item The compressor obtains
 $(x_{n,c})_c=\mathsf{TwoForestSplit}_{n_0,u}(\mathit{DB}^{*})$.  For each
 sampled $c$, set
 \[
 C_{n,c}:=
 \hyperref[alg:additive]{\textsc{AdditiveEC}}
 (\mathsf{encode},(\ell_c,x_{n,c}))
 \]
 and record $C_{n,c}[I_c]$.  Return the concatenation of these vectors in
 public chunk order.
 \item The extractor parses the advice into vectors $(v_c)_c$.  For sampled
 $c$, set
 \[
 x'_c:=
 \hyperref[alg:additive]{\textsc{AdditiveEC}}(\mathsf{decode},
 (\ell_c,\mathsf{Fill}_c(Y_c,I_c,v_c))).
 \]
 For mini $c$, set
 $x'_c:=\hyperref[alg:additive]{\textsc{AdditiveEC}}
 (\mathsf{decode},(\ell_c,Y_c))$ for its partial redundancy-two word.  For
 tiny $c$, set $x'_c$ to its raw message if every
 coordinate is present and to $\bot$ otherwise.  Return
 $\mathsf{TwoForestJoin}_{n_0,u}((x'_c)_c)$ if every $x'_c\ne\bot$, and return
 $\bot$ otherwise.
\end{enumerate}
\end{algbox}

The following lemma combines the deterministic decoding condition with the
probability and advice accounting established above.

\begin{lemma}[Paired recovery format, decoding, and resources]
\label{lem:recovery}
Under all recovery hypotheses of \Cref{thm:main}, the two entries of
\hyperref[alg:recovery]{\textsc{TwoForestRecovery}} are total.  On paired
inputs, the advice length is exactly
\[
 B\sum_{c\text{ sampled}}d_c
 \le B\max\{n-(1-\epsilon)\mu S,0\},
\tag{26}\label{eq:advice}
\]
and the advice contains no indices.  The extractor runs in time
quasi-polynomial in $(\lambda,t_{\mathcal A})$ and satisfies
\[
 \Pr[\mathsf{Pass}\wedge
 (\mathit{DB}_{\rm ext}\ne\mathit{DB}^{*})]
 \le\operatorname{negl}(\lambda).
\tag{27}\label{eq:recovery-probability}
\]
\end{lemma}

\begin{proof}
All parsers, searches, rewinds, completion maps, decoders, and malformed-input
branches have finite public bounds, so both entries are total.  Fix the good
events in \Cref{lem:effective-support}.  For a sampled chunk, paired advice
supplies the correct codeword values on $I_c$, disjoint from $P_c$.  The filled
word has $|P_c|+d_c$ supported entries and exactly $e_c$ errors.  By
\Cref{eq:effective-support},
\[
 |P_c|+d_c-2e_c\ge c_c+(n_c-c_c)=n_c.
\]
The decoder guarantee in \Cref{lem:additive-code} therefore returns $x_{n,c}$.
The mini inequality and exact tiny opening in \Cref{lem:effective-support}
give the same conclusion for those chunks.  \Cref{lem:partition} then makes
the final join equal $\mathit{DB}^{*}$.

For the advice count, \Cref{eq:sampling-properties} gives
\[
 c_c\ge n_c\min\left\{1,\frac{\vartheta\mu S}{n}\right\}
\]
on every sampled chunk, while mini and tiny chunks use no advice.  Since the
chunk lengths sum to $n$,
\[
 \sum_{c\text{ sampled}}d_c
 \le\max\{n-\vartheta\mu S,0\}
 \le\max\{n-(1-\overline\epsilon)\mu S,0\}
 \le\max\{n-(1-\epsilon)\mu S,0\},
\]
where $\vartheta=1-13\overline\epsilon/40\ge1-\overline\epsilon$ and
$\overline\epsilon\le\epsilon$.  This proves \Cref{eq:advice}, including
zero advice in the saturated regime.

There are at most $2(1+\lfloor\log_2N\rfloor)$ chunks.  Combining
\Cref{def:local-extraction,lem:challenge-sampling,lem:spine,lem:deferred}
bounds the bad event jointly
with $\mathsf{Pass}$ by
\[
 2^{-\lambda-1}+\frac1{K_{\rm sp}+1}
 +K_{\rm sp}q_{\rm loc}
   \left(\frac{8N}{T_{\rm ext}+1}+(1-\eta)^\kappa\right)
 +\frac{K_{\rm sp}\mu}{T_{\rm ext}+1}
 +\nu_{\rm bind}(\lambda)
 +\sum_{c\text{ sampled}}\nu_{{\rm exp},c}(\lambda).
\tag{28}\label{eq:total-failure}
\]
Every term is negligible by \Cref{eq:spine-balance},
\Cref{eq:exposure-failure}, polynomiality of $N,\mu$, and
quasi-polynomial binding.  Outside this event the deterministic decoding above
succeeds, proving \Cref{eq:recovery-probability}.

The accepting-spine and rewind work is quasi-polynomial by \Cref{lem:spine}.
Completion scans $O(Rn_c)$ coordinates per chunk, and
\Cref{lem:additive-code} bounds decoding by
\[
 B^{O(1)}\sum_{c\in\mathcal K(u)}(Rn_c)^4
\]
bit operations.  The field-size, public-size, and polynomial-redundancy
conditions make this polynomial in the displayed parameters; moreover
$B\le t_{\mathcal A}$ because the experiment writes a nonempty $B$-bit-block
database.  Hence the total extractor time is quasi-polynomial in
$(\lambda,t_{\mathcal A})$.  Deterministic paired equality from
\Cref{lem:spine,lem:effective-support} shows that no support metadata or
transcript is needed by the extractor.
\end{proof}

We can now substitute the rounded public slack and finish the theorem.

\begin{proof}[Proof of \Cref{thm:main}]
The definition of $\overline\epsilon$ gives
$\epsilon/2<\overline\epsilon\le\epsilon$.  Consequently
\[
 R=\Theta(1/\overline\epsilon)=\Theta(1/\epsilon),
 \qquad
 \zeta_0=2^{-(k_{\rm rec}+3)}\in(\epsilon/16,\epsilon/8].
\]
The dyadic parameter has a canonical encoding of
$O(1+\log(1/\epsilon))$ bits.  Since
$R\widehat N\le2^B$, one has $B=\Omega(\log N+\log(1/\epsilon))$, so
$z_0=O(B)$.  Substitution into \Cref{eq:aggregate-update} proves
$O(BU\log N/\epsilon)$ simplified publisher work for every $U$, and the bit
work, storage, and tree-depth conclusions follow from \Cref{lem:publisher}.
The proof counts only suffix carries and therefore introduces no relation
between $n_0$ and $U$.  The same lemma proves honest correctness and, under
the additional recovery block-size condition, the node-state and
externally-materialized output bounds.

For recovery, the extra hypothesis implies
$\overline\epsilon\kappa=\omega(\log\lambda)$.  Also
$R\ge12\mathrm e/\eta$, $\vartheta>1/2$, and all recovery-validity and
efficiency conditions needed by \Cref{lem:recovery} hold uniformly at every
legal length.  That lemma gives the exact paired interfaces, advice bound,
quasi-polynomial extractor, and negligible joint error asserted in the
theorem.  Its construction and proof use pairwise distinct identities and make
no replication claim, which establishes the stated scope.
\end{proof}

\bibliographystyle{alpha}
\bibliography{references}

\end{document}
